% Editable source reconstructed from cycles-through-1-qn1.pdf.
% Display formulas follow the printed pages. Four inline powers in that PDF
% are pandoc brace errors (the superscript swallowed the next token). They
% are written here as the identities the proofs use:
%   Lemma 11:        2^{S_u}/q^{u+1}  and  x_i 2^{S_i}/2^{ih/\ell}
%   Theorem 2:       2^m (2^{r/\ell} - (\ell q+1)/(\ell 2^m))
%   Theorem 3:      (s 2^j)^{-1-\varepsilon}
%   Proposition B.12: 2^{k_1}/2^{k_2}
\documentclass[11pt]{article}

\usepackage[margin=1in]{geometry}
\usepackage{lmodern}
\usepackage{microtype}
\usepackage{amsmath,amssymb,amsthm}
\usepackage{amsopn}
\usepackage{unicode-math}
\usepackage{amstext}
% A blank line in the source opens a line of space on the page.
% Section headings are adjusted so they do not pick up a second skip.
\usepackage[skip=\baselineskip]{parskip}
% Long inline formulas may break at a relation or a binary operator,
% instead of stretching the spaces on that line.
\relpenalty=100
\binoppenalty=200
% Reference links only. Section and statement cross-references stay black.
\usepackage[bookmarks=false,colorlinks,linkcolor=black,citecolor=black,urlcolor=blue]{hyperref}
\newcommand{\doilink}[1]{\href{https://doi.org/#1}{\nolinkurl{doi:#1}}}

\newenvironment{named}[1]{%
  \par\addvspace{\medskipamount}%
  \pagebreak[3]%
  \noindent{\bfseries #1}\space\ignorespaces
  \stmtoc{#1}%
}{\par\addvspace{\medskipamount}}

\newcommand{\topic}[1]{\par\noindent\textbf{#1}}

\setcounter{tocdepth}{2}
\ExplSyntaxOn
\prop_new:N \g_qn_toc_total_prop
\prop_new:N \g_qn_toc_seen_prop
\bool_new:N \g_qn_toc_appendix_bool
\AddToHook { cmd/appendix/before } { \bool_gset_true:N \g_qn_toc_appendix_bool }
\cs_new_protected:Npn \qnTocCount #1 #2
  { \prop_gput:Nnn \g_qn_toc_total_prop {#1} {#2} }
\cs_new_protected:Npn \qn_stmtoc:n #1
  {
    \tl_set:Nn \l_tmpa_tl {#1}
    \bool_set_false:N \l_tmpa_bool
    \tl_if_in:NnT \l_tmpa_tl { Outline } { \bool_set_true:N \l_tmpa_bool }
    \tl_if_in:NnT \l_tmpa_tl { Bauer } { \bool_set_true:N \l_tmpa_bool }
    \tl_if_in:NnT \l_tmpa_tl { Ridout } { \bool_set_true:N \l_tmpa_bool }
    \bool_if:NF \l_tmpa_bool
      {
        \regex_match:nVTF { \A (Conjecture|Theorem|Proposition)\b | conjecture } \l_tmpa_tl
          {
            \bool_if:NTF \g_qn_toc_appendix_bool
              {
                \regex_match:nVTF { \A Proposition \b } \l_tmpa_tl
                  { }
                  { \qn_stmtoc_write: }
              }
              { \qn_stmtoc_write: }
          }
          { }
      }
  }
\cs_new_protected:Npn \qn_stmtoc_write:
  {
    \tl_set_eq:NN \l_tmpb_tl \l_tmpa_tl
    \regex_replace_once:nnN { \.\Z } { } \l_tmpb_tl
    \int_set:Nn \l_tmpa_int { 1 }
    \prop_get:NVNT \g_qn_toc_seen_prop \l_tmpb_tl \l_tmpa_tl
      { \int_set:Nn \l_tmpa_int { \l_tmpa_tl + 1 } }
    \tl_set:Nx \l_tmpa_tl { \int_use:N \l_tmpa_int }
    \exp_args:NNVV \prop_gput:Nnn \g_qn_toc_seen_prop \l_tmpb_tl \l_tmpa_tl
    \use:e
      {
        \exp_not:N \qnTocWrite
        { \exp_not:V \l_tmpb_tl }
        { \exp_not:V \l_tmpa_tl }
      }
    \prop_get:NVNT \g_qn_toc_total_prop \l_tmpb_tl \l_tmpa_tl
      {
        \int_compare:nNnT { \l_tmpa_int } = { \l_tmpa_tl }
          {
            \phantomsection
            \exp_args:Nnx \addcontentsline {toc} {statement}
              { \exp_not:V \l_tmpb_tl }
          }
      }
  }
\newcommand{\stmtoc}[1]{ \qn_stmtoc:n {#1} }
\ExplSyntaxOff
\makeatletter
\newcommand{\qnTocWrite}[2]{%
  \immediate\write\@auxout{\string\qnTocCount{\unexpanded{#1}}{#2}}%
}
\newcommand*\l@statement[2]{%
  \ifnum \c@tocdepth > 1\relax
    \addpenalty\@secpenalty
    \begingroup
      \parindent \z@
      \leftskip 2.6em
      \rightskip \@pnumwidth
      \parfillskip -\@pnumwidth
      \leavevmode
      #1\nobreak\hfil\nobreak
      \hb@xt@\@pnumwidth{\hss #2}\par
    \endgroup
  \fi
}
\makeatother

\newcounter{stmtref}
% The headings above are written out by hand, so a label would otherwise
% capture the section number. This records the printed statement number.
\newcommand{\markstmt}[2]{%
  \setcounter{stmtref}{#1}%
  \addtocounter{stmtref}{-1}%
  \refstepcounter{stmtref}%
  \label{#2}%
}

\title{I. Cycles through 1 for the $qn+1$ map: no cycles of length $3$, and constraints on longer cycles}
\author{Trevor Hackman}
\date{}

\begin{document}
\maketitle

\begin{abstract}
\noindent
For the $qn+1$ map with $q$ an odd integer at least $3$, we prove that no cycle with exactly three distinct odd numbers has starting value $1$.
This settles the case of length $3$ in a conjecture of Santos (2020) that a cycle has starting value $1$ only when $q = 5$ or when $q$ is a Mersenne number.
We found that for every such $q$ at least $2^{22001}$, a lower bound and an upper bound cannot hold at the same time.
Bauer and Bennett (2018) give the exponent $\mu \approx 1.61487$, which lies just below the golden ratio, the threshold at which these bounds conflict.
Every smaller $q$ is ruled out by a direct computation.
\end{abstract}

\begingroup
\setlength{\parskip}{0pt}
\small
\tableofcontents
\endgroup

\section{Introduction}
\label{sec:introduction}

This paper studies cycles of the $qn+1$ map with starting value $1$.
It is followed by the broader paper~\cite{hackman-ii}, which studies cycles of the same map for any starting value.

Fix an odd integer $q \ge 3$.
For an odd positive integer $n$, let $v$ be the largest integer such that $2^v$ divides $qn+1$.
The map $C$ sends $n$ to
\[
C(n) = \frac{qn+1}{2^v}.
\]
The orbit of $n$ is the sequence $n$, $C(n)$, $C(C(n))$, \ldots.
When $q = 3$, $C$ is the odd-to-odd form of the Collatz map, also called the Syracuse function, and the Collatz conjecture, posed by Lothar Collatz in 1937 and still open, asserts that every orbit eventually reaches $1$.
For $q = 3$, heuristic arguments suggest that almost all orbits eventually reach $1$.
However, for $q \ge 5$, heuristic arguments suggest that almost all orbits grow without bound and that cycles are rare, although no orbit has been proved to diverge.

We ask for which $q$ a cycle of $C$ has starting value $1$.
We classify the answers by the length $\ell$ of the cycle with starting value $1$, that is, the number of distinct odd numbers on it.

For $\ell = 1$ the cycle is $1 \to 1$, which happens exactly when $q$ is a Mersenne number (\hyperref[prop:a1]{Proposition~A.1}).
For $\ell = 2$, the only $q$ is $5$, with $1 \to 3 \to 1$ (\hyperref[prop:a2]{Proposition~A.2}).
Besides $1 \to 1$ when $q$ is a Mersenne number, the only known cycle with starting value $1$ is $1 \to 3 \to 1$ for $q = 5$.

\begin{named}{Conjecture 1 (Santos (2020)).}
A cycle of $C$ has starting value $1$ only when $q = 5$ or when $q$ is a Mersenne number.
\end{named}

\begin{named}{Theorem 1.}
No odd $q \ge 3$ has a cycle of $C$ of length $3$ with starting value $1$.
\end{named}

\begin{named}{Outline of Theorem 2 (golden-ratio criterion).}
Fix $\ell \ge 2$, and suppose that length $\ell$ has a restricted irrationality exponent below the golden ratio.
Then any $q$ that has a cycle of length $\ell$ with starting value $1$ lies below a bound we can compute.
Checking every candidate up to that bound finds every such $q$, or shows that none exist.
The precise statement is in Section~\ref{sec:longer}.
\end{named}

\begin{named}{Theorem 3 (finiteness for starting value $1$).}
For each $\ell \ge 2$, only finitely many odd $q$ have a cycle of $C$ of length $\ell$ with starting value $1$.
\end{named}

Theorem~3 is ineffective: it rests on the $S$-unit form of Ridout's theorem in Section~\ref{sec:ineffective}, and gives no bound on $q$.
(A finiteness result is called effective if its proof gives an explicit bound on any exception, so that a finite computation can find them all, and ineffective if it gives no such bound.)
Theorem~1 is effective and is completed by computation.
The proof of Theorem~3, in Section~\ref{sec:ineffective}, uses only the $S$-unit form of Ridout's theorem and Theorem~2.
Cycles with any starting value are treated in the companion paper~\cite{hackman-ii}.
For a cycle of length $2$ or $3$ with a constant odd starting value, that paper gives a bound on $q$ depending only on the element and the length (Theorem~3), and for Corollaries~12 to~14 the same bound is Corollary~14 of that paper.

\section{Definitions and notation}
\label{sec:definitions}

Throughout, $q \ge 3$ is an odd integer and $C$ is the map defined in Section~1.

\topic{Mersenne numbers.}
Let $e \ge 2$ be an integer.
A Mersenne number is the integer $2^e - 1$.

\topic{Divides.}
The symbol $\mid$ means divides, and the symbol $\nmid$ means does not divide.
Let $a$ and $b$ be integers, with $a \ne 0$.
Write $a \mid b$ when there exists an integer $k$ such that $b = ak$.
Write $a \nmid b$ when no such integer $k$ exists.

\topic{The $2$-adic valuation.}
For a nonzero integer $n$, $v_2(n)$ is the largest integer $v$ such that $2^v$ divides $n$.

\topic{Cycles.}
A cycle of $C$ of length $\ell$ consists of $\ell$ distinct odd positive integers.
The starting value is the smallest element $x_0$.
Applying $C$ repeatedly, starting at $x_0$, produces
\[
x_0 \to x_1 \to \cdots \to x_{\ell-1} \to x_0,
\]
where the indices are read modulo $\ell$.

\topic{Starting value $1$.}
A cycle with starting value $1$ is a cycle of $C$ of length $\ell$ with $x_0 = 1$.
It is written
\[
1 = x_0 \to x_1 \to \cdots \to x_{\ell-1} \to 1.
\]

\topic{Halvings.}
The halving counts $k_1, \ldots, k_\ell$ of the cycle are defined by
\[
q x_{i-1} + 1 = 2^{k_i} x_i
\qquad (1 \le i \le \ell,\ x_\ell = x_0),
\]
so $k_i = v_2(q \cdot x_{i-1} + 1)$.
The total number of halvings is $h = k_1 + \cdots + k_\ell$, and $S_i = k_1 + \cdots + k_i$.

\topic{The $\delta$ window.}
Let $\ell \ge 2$ be an integer.
On a cycle of length $\ell$ with total number of halvings $h$, the number $\delta$ is the fractional part of $2^{h/\ell}$.
The integer $q$ equals $\lfloor 2^{h/\ell} \rfloor$, so $2^{h/\ell} = q + \delta$ and $0 < \delta < 1$.
The $\delta$ window is a lower bound and an upper bound on $\delta$.

\topic{Near misses.}
For an element $x_i$ other than $1$ on a cycle of length $\ell$ with starting value $1$, set $n_i = ih - \ell S_i$.
The near miss of $x_i$ is $2^{n_i/\ell}$.
The relative error $\rho_i$ is defined by
\[
x_i = 2^{n_i/\ell}(1 + \rho_i).
\]

\topic{The restricted irrationality exponent.}
Fix an integer $\ell \ge 2$.
A real number $\mu$ is a restricted irrationality exponent for length $\ell$ when, for each odd positive integer $x$, there is a positive real number $c(x)$ and there is a nonnegative integer $j_0(x)$ such that, for every integer $r$ with $1 \le r \le \ell - 1$, every integer $p$, and every integer $j \ge j_0(x)$,
\[
\left| 2^{r/\ell} - \frac{p}{\ell x \cdot 2^j} \right| > c(x) \cdot 2^{-\mu j}.
\]

\topic{The golden ratio.}
The golden ratio is $\varphi = (1 + \sqrt{5})/2$.

\topic{The order of $2$.}
For an odd positive integer $m$, the order of $2$ modulo $m$ is the smallest positive integer $d$ such that $2^d \equiv 1 \pmod{m}$.

\topic{Big $O$.}
Big $O$ is the Bachmann--Landau symbol $O$.
Let $f$ and $g$ be real-valued functions of a real variable $t$, with $g(t) \ne 0$ for every large $t$.
Write $f(t) = O(g(t))$ when there exists a constant real number $C > 0$, independent of $t$, and there exists a constant real number $t_0$, such that $|f(t)| \le C\,|g(t)|$ for every $t \ge t_0$.
The expression $O(g(t))$ stands for a real number of absolute value at most $C\,|g(t)|$.
An equation $A = B + O(g(t))$ means $A - B = O(g(t))$.

\topic{Vinogradov's symbols.}
Vinogradov's symbols are $\ll$ and $\gg$.
Let $f$ and $g$ be real-valued functions of a real variable $t$.
Write $f \ll g$ when $f = O(g)$.
Write $f \gg g$ when $g \ll f$.

\topic{Hardy's symbol.}
Hardy's symbol is $\asymp$.
Write $f \asymp g$ when $f \ll g$ and $g \ll f$.

\topic{Less than or similar to.}
The symbol $\lesssim$ means less than or similar to, and the symbol $\gtrsim$ means greater than or similar to.
Let $f$ and $g$ be positive real numbers.
Write $f \lesssim g$ when $f$ is at most a constant real number times $g$, and that constant is not named.
Write $f \gtrsim g$ when $g \lesssim f$.

\section{The structure of a cycle of length $3$}
\label{sec:structure}

\topic{The three steps of length $3$.}
On a cycle of length $3$ with starting value $1$, write $x$ for $x_1$ and $y$ for $x_2$.
The cycle is $1 \to x \to y \to 1$.
The three steps are $q + 1 = 2^{k_1} x$, $qx + 1 = 2^{k_2} y$, and $qy + 1 = 2^{k_3}$.

\begin{named}{Lemma 1 (lower bound for $k_3$).}
$k_3 \ge k_1 + 3$.
\end{named}

\begin{proof}
The odd elements other than $1$ are at least $3$, so $x \ge 3$ and $y \ge 3$.
The first step is $q + 1 = 2^{k_1} x$, and therefore $q + 1 \ge 3 \cdot 2^{k_1}$, that is $q \ge 3 \cdot 2^{k_1} - 1$.
The third step is $2^{k_3} = qy + 1$, so
\[
2^{k_3} \ge (3 \cdot 2^{k_1} - 1) \cdot 3 + 1 = 9 \cdot 2^{k_1} - 2.
\]
Since $2^{k_1} \ge 2$, we have $9 \cdot 2^{k_1} - 2 \ge 8 \cdot 2^{k_1} = 2^{k_1 + 3}$, and hence $k_3 \ge k_1 + 3$.
\end{proof}

\begin{named}{Lemma 2 (cycle equation).}
\[
2^h = q^3 + q^2 + 2^{k_1} q + 2^{k_1 + k_2}.
\]
\end{named}

\begin{proof}
Since $h = k_1 + k_2 + k_3$,
\[
2^h = 2^{k_1 + k_2} \cdot 2^{k_3}.
\]
The third step is $qy + 1 = 2^{k_3}$, so
\[
2^h = 2^{k_1 + k_2}(qy + 1) = 2^{k_1}\bigl(q \cdot 2^{k_2} y + 2^{k_2}\bigr).
\]
The second step is $qx + 1 = 2^{k_2} y$, and therefore
\[
2^h = 2^{k_1}\bigl(q(qx + 1) + 2^{k_2}\bigr) = 2^{k_1}(q^2 x + q + 2^{k_2}) = q^2(2^{k_1} x) + 2^{k_1} q + 2^{k_1 + k_2}.
\]
The first step is $q + 1 = 2^{k_1} x$, so
\[
2^h = q^2(q + 1) + 2^{k_1} q + 2^{k_1 + k_2} = q^3 + q^2 + 2^{k_1} q + 2^{k_1 + k_2}.
\]
\end{proof}

\begin{named}{Lemma 3 (between consecutive cubes).}
\markstmt{3}{lem:cube-window}
$q^3 < 2^h < (q+1)^3$.
Hence $3$ does not divide $h$ and $q = \lfloor 2^{h/3} \rfloor$.
\end{named}

\begin{proof}
Lemma~2 gives $2^h = q^3 + q^2 + 2^{k_1} q + 2^{k_1 + k_2}$.
The terms $q^2$, $2^{k_1} q$, and $2^{k_1 + k_2}$ are positive, so $2^h > q^3$.

Since $h = k_1 + k_2 + k_3$,
\[
2^h = 2^{k_1} \cdot 2^{k_2} \cdot 2^{k_3}.
\]
The three steps give $2^{k_1} = (q+1)/x$, $2^{k_2} = (qx+1)/y$, and $2^{k_3} = qy + 1$, so
\[
2^h = \frac{q+1}{x} \cdot \frac{qx+1}{y} \cdot (qy + 1) = (q+1) \cdot \frac{qx+1}{x} \cdot \frac{qy+1}{y}.
\]
Since $(qx+1)/x = q + 1/x$ and $(qy+1)/y = q + 1/y$,
\[
2^h = (q+1)\left(q + \frac{1}{x}\right)\left(q + \frac{1}{y}\right).
\]
The elements $x$ and $y$ are distinct from the starting value $1$, so $x > 1$ and $y > 1$.
Hence $1/x < 1$ and $1/y < 1$, so $q + 1/x < q + 1$ and $q + 1/y < q + 1$, and $2^h < (q+1)^3$.

A power of $2$ strictly between consecutive cubes is not a cube, so $3 \nmid h$.
Since $q^3 < 2^h < (q+1)^3$, $q = \lfloor 2^{h/3} \rfloor$.
\end{proof}

\begin{named}{Lemma 4 ($2$-adic valuation of $x-1$).}
Let $t$ be the integer $v_2(x-1)$, and let $s$ be the odd positive integer $(x-1)/2^t$.
Then either $t = 2k_1$, or $t < k_1$ and $k_1 = t + v_2(s+1)$, or $t = k_1$.
\end{named}

\begin{proof}
The first step is $q + 1 = 2^{k_1} x$, so $q = 2^{k_1} x - 1$.

Substitute $q = 2^{k_1} x - 1$ into the second step $2^{k_2} y = qx + 1$:
\[
2^{k_2} y = (2^{k_1} x - 1)x + 1 = 2^{k_1} x^2 - x + 1 = 2^{k_1} x^2 - (x - 1).
\]

Suppose $t > k_1$.
Since $x$ is odd, $v_2(2^{k_1} x^2) = k_1$, while $v_2(x-1) = t > k_1$, so
\[
2^{k_1} x^2 - (x-1) = 2^{k_1}\bigl(x^2 - 2^{t-k_1}s\bigr),
\]
and $x^2 - 2^{t-k_1}s$ is odd, since $x^2$ is odd and $2^{t-k_1}s$ is even.
Thus the $2$-adic valuation of $2^{k_2} y$ is $k_1$. Since $y$ is odd, $k_2 = k_1$, so $2^{k_1} y = 2^{k_1} x^2 - (x-1)$.

Let $w = (x-1)/2^{k_1}$. Then $w$ is even, since $t > k_1$, and $x = 1 + 2^{k_1} w$.
Also $y = x^2 - w$, so
\[
y = (1 + 2^{k_1} w)^2 - w = 1 + w(2^{k_1+1} - 1) + 2^{2k_1} w^2.
\]
The first step gives $q = 2^{k_1} x - 1 = 2^{k_1} + 2^{2k_1} w - 1$.
Multiplying these expressions for $q$ and $y$ and adding $1$ gives
\[
qy + 1 = 2^{k_1} + w(1 - 3 \cdot 2^{k_1}) + 2^{2k_1} w(3 - 2w) + 3 \cdot 2^{3k_1} w^2 + 2^{4k_1} w^3.
\]
The terms $2^{2k_1} w(3 - 2w)$, $3 \cdot 2^{3k_1} w^2$, and $2^{4k_1} w^3$ are divisible by $2^{2k_1 + v_2(w)}$, and $1 - 3 \cdot 2^{k_1}$ is odd, so $v_2\bigl(w(1 - 3 \cdot 2^{k_1})\bigr) = v_2(w)$. If $v_2(w) \ne k_1$, then $v_2(qy + 1) = \min(k_1, v_2(w)) \le k_1$, contradicting Lemma~1.
Hence $v_2(w) = k_1$, that is $t = 2k_1$.

Suppose $t < k_1$.
Then $2^{k_2} y = 2^{k_1} x^2 - 2^t s = 2^t\bigl(2^{k_1-t}x^2 - s\bigr)$.
The factor $2^{k_1-t}x^2 - s$ is odd, since $2^{k_1-t}x^2$ is even and $s$ is odd, so $k_2 = t$ and $y = 2^{k_1-t}x^2 - s$.
Since $x = 1 + 2^t s$, we have $q = 2^{k_1} x - 1 = (2^{k_1} - 1) + 2^{k_1 + t} s$, and therefore
\[
qy + 1 = (2^{k_1}-1)y + 2^{k_1 + t} s y + 1.
\]
Substitute $y = 2^{k_1-t}x^2 - s$:
\[
qy + 1 = (2^{k_1}-1)\, 2^{k_1-t} x^2 + (1+s) - 2^{k_1} s + 2^{k_1 + t} s y.
\]
Now $2^{k_1} s$ and $2^{k_1 + t} s y$ are divisible by $2^{k_1}$, and $v_2\bigl((2^{k_1}-1)\, 2^{k_1-t} x^2\bigr) = k_1 - t$, since $2^{k_1}-1$ and $x$ are odd. If $v_2(s+1) \ne k_1 - t$, then $v_2(qy + 1) < k_1$, contradicting Lemma~1.
Hence $v_2(s+1) = k_1 - t$, that is $k_1 = t + v_2(s+1)$.

If $t = k_1$, that is the remaining case.
\end{proof}

\begin{named}{Lemma 5 (lower bound for $x$).}
$2^{k_1} \le 2(x-1)$ and $x > \sqrt{q/2}$.
\end{named}

\begin{proof}
If $t = 2k_1$, then $2^{2k_1} = 2^t \le x - 1$, so $2^{k_1} \le x - 1 \le 2(x-1)$.
If $t < k_1$ and $k_1 = t + v_2(s+1)$, then $2^{k_1} = 2^t \cdot 2^{v_2(s+1)} \le 2^t(s+1) = (x-1) + 2^t \le 2(x-1)$.
If $t = k_1$, then $2^{k_1} = 2^t \le x - 1 \le 2(x-1)$.
In every case $2^{k_1} \le 2(x-1)$.
The first step is $q + 1 = 2^{k_1} x$, so
\[
q + 1 \le 2x(x-1) = 2x^2 - 2x.
\]
Hence $q < 2x^2$, so $x > \sqrt{q/2}$.
\end{proof}

\begin{named}{Lemma 6.}
Let $\ell \ge 2$.
Let $i$ be an integer with $1 \le i \le \ell - 1$.
On a cycle of length $\ell$ with starting value $1$,
\[
2^{k_{i+1}} x_{i-1}\, x_{i+1} - 2^{k_i} x_i^2 = x_{i-1} - x_i.
\]
Also $|x_{i-1} - x_i| \ge 2^{\min(k_i, k_{i+1})}$.
If $k_i \le k_{i+1}$, then $x_i > q$ when $x_i < x_{i-1}$, and $x_i^2 > q\, x_{i-1}$ when $x_i > x_{i-1}$.
\end{named}

\begin{proof}
Since $q x_i + 1 = 2^{k_{i+1}} x_{i+1}$, we have $q = (2^{k_{i+1}} x_{i+1} - 1)/x_i$.
Substitute $q = (2^{k_{i+1}} x_{i+1} - 1)/x_i$ into $q x_{i-1} + 1 = 2^{k_i} x_i$ and multiply by $x_i$:
\[
2^{k_{i+1}} x_{i-1}\, x_{i+1} - x_{i-1} + x_i = 2^{k_i} x_i^2,
\]
and therefore $2^{k_{i+1}} x_{i-1}\, x_{i+1} - 2^{k_i} x_i^2 = x_{i-1} - x_i$.

The integers $x_{i-1}$, $x_i$, and $x_{i+1}$ are odd, so $x_{i-1} x_{i+1}$ and $x_i^2$ are odd.
The $2$-adic valuation of $2^{k_{i+1}} x_{i-1} x_{i+1}$ is $k_{i+1}$, and the $2$-adic valuation of $2^{k_i} x_i^2$ is $k_i$.
If $k_i \ne k_{i+1}$, the $2$-adic valuation of $2^{k_{i+1}} x_{i-1} x_{i+1} - 2^{k_i} x_i^2$ is $\min(k_i, k_{i+1})$.
If $k_i = k_{i+1}$, then $x_{i-1} x_{i+1} - x_i^2$ is even, since it is a difference of two odd integers, so the $2$-adic valuation of $2^{k_i}(x_{i-1} x_{i+1} - x_i^2)$ is at least $k_i$.
In either case $2^{\min(k_i, k_{i+1})}$ divides $x_{i-1} - x_i$.
The elements of the cycle are distinct, so $x_{i-1} \ne x_i$ and $|x_{i-1} - x_i| \ge 2^{\min(k_i, k_{i+1})}$.

Suppose $k_i \le k_{i+1}$.
Then $|x_{i-1} - x_i| \ge 2^{k_i} = (q x_{i-1} + 1)/x_i$.
If $x_i > x_{i-1}$, then $x_i(x_i - x_{i-1}) \ge q x_{i-1} + 1$.
Since $x_{i-1} \ge 1$, we have $x_i^2 > x_i(x_i - x_{i-1})$, and therefore $x_i^2 > q x_{i-1}$.
If $x_i < x_{i-1}$ and $x_i \le q$, then $x_{i-1} x_i \le q x_{i-1}$, so
\[
x_i(x_{i-1} - x_i) = x_{i-1} x_i - x_i^2 \le q x_{i-1} - x_i^2 < q x_{i-1} + 1.
\]
This contradicts $x_i(x_{i-1} - x_i) \ge q x_{i-1} + 1$.
Hence $x_i > q$.
\end{proof}

\begin{named}{Lemma 7.}
$y > x$ and $y^2 > qx$.
In particular $2^{k_2} < q + 1$, so $2^{k_1 + k_2} \le 2^{k_1} q$.
\end{named}

\begin{proof}
Apply Lemma~6 with $i = 2$. Then $x_1 = x$, $x_2 = y$, and $x_3 = 1$.
Suppose $k_3 < k_2$. Lemma~6 gives $|x - y| \ge 2^{k_3}$.
The third step is $qy + 1 = 2^{k_3}$, so $|x - y| \ge qy + 1$.
If $y > x$, then $y - x \ge qy + 1$, and $y(1-q) \ge x + 1 > 0$.
Here $q \ge 3$ and $y > 0$, so $y(1-q) < 0$, a contradiction.
Thus $x > y$, and $x - y \ge qy + 1$, so $x \ge (q+1)y + 1 > qy$.
The element $y$ is an odd integer at least $3$, so $x > 3q$.
The first step is $q + 1 = 2^{k_1} x$ with $k_1 \ge 1$, so $x \le (q+1)/2$.
Since $q \ge 3$, we have $3q > (q+1)/2$, and therefore $x > (q+1)/2 \ge x$, a contradiction.
Thus $k_2 \le k_3$.

If $y < x$, Lemma~6 gives $y > q$.
Also $x = (q+1)/2^{k_1} \le (q+1)/2 < q$, so $y > q > x$, a contradiction.
The elements of the cycle are distinct, so $y \ne x$.
Hence $y > x$, and Lemma~6 gives $y^2 > qx$.
Since $y > x$,
\[
2^{k_2} = \frac{qx + 1}{y} < \frac{qx + 1}{x} = q + \frac{1}{x} < q + 1.
\]
The integer $2^{k_2}$ is at most $q$, so $2^{k_1 + k_2} \le 2^{k_1} q$.
\end{proof}

\section{The $\delta$ window and the near miss for $x$}
\label{sec:window}

\begin{named}{Lemma 8 (the $\delta$ window).}
\markstmt{8}{lem:window}
Let $N = 27 \cdot 2^h - (3q+1)^3$.
Then
\[
N = 9q(3 \cdot 2^{k_1}-1) + 27 \cdot 2^{k_1 + k_2} - 1,
\qquad
0 < N < 54 \cdot 2^{k_1} q,
\qquad
0 < \delta - \frac{1}{3} \le \frac{N}{81q^2} < \frac{2^{k_1 + 1}}{3q}.
\]
\end{named}

\begin{proof}
Lemma~2 gives $2^h = q^3 + q^2 + 2^{k_1} q + 2^{k_1 + k_2}$.
Expanding $(3q+1)^3 = 27q^3 + 27q^2 + 9q + 1$ and subtracting from $27 \cdot 2^h$ gives
\[
N = 9q(3 \cdot 2^{k_1}-1) + 27 \cdot 2^{k_1 + k_2} - 1.
\]
Since $2^{k_1} \ge 2$, the formula gives $N \ge 9q(3 \cdot 2 - 1) - 1 = 45q - 1 > 0$.
Lemma~7 gives $2^{k_1 + k_2} \le 2^{k_1} q$, so
\[
N < 9q(3 \cdot 2^{k_1}) + 27 \cdot 2^{k_1} q = 54 \cdot 2^{k_1} q.
\]
Since $2^{h/3} = q + \delta$ and $(3q+1)^3 = 27(q + 1/3)^3$,
\[
N = 27\bigl[(q+\delta)^3 - (q + 1/3)^3\bigr].
\]
The factor $(q+\delta)^3 - (q + 1/3)^3$ equals $(\delta - 1/3)$ times $(q+\delta)^2 + (q+\delta)(q + 1/3) + (q + 1/3)^2$.
Each of those three terms is greater than $q^2$, so their sum is greater than $3q^2$.
Thus $N > 0$ gives $\delta > 1/3$, and $N \ge 81q^2(\delta - 1/3)$, so $\delta - 1/3 \le N/(81q^2)$.
Also
\[
\frac{N}{81q^2} < \frac{54 \cdot 2^{k_1} q}{81q^2} = \frac{2^{k_1+1}}{3q}.
\]
\end{proof}

\begin{named}{Lemma 9 (near miss for $x$).}
\markstmt{9}{lem:near-x}
Let $f = h - 3k_1$.
Then $3 \nmid f$, and
\[
\bigl|x^3 - 2^f\bigr| \le \frac{3.2\, x^2}{2^{k_1}},
\qquad
\bigl|x - 2^{f/3}\bigr| \le \frac{1.6}{2^{k_1}}.
\]
\end{named}

\begin{proof}
Since $q + 1 = 2^{k_1} x$, we have $2^{3k_1} x^3 = (q+1)^3$.

Let $E = 2q^2 + (3 - 2^{k_1})q + 1 - 2^{k_1 + k_2}$.
Lemma~2 gives $2^h = q^3 + q^2 + 2^{k_1} q + 2^{k_1 + k_2}$, so
\[
2^{3k_1} x^3 - 2^h = (q+1)^3 - 2^h = 2q^2 + (3 - 2^{k_1})q + 1 - 2^{k_1 + k_2} = E.
\]
Since $2^{k_1} \ge 2$, either $2^{k_1} = 2$ and $|3 - 2| = 1$, or $2^{k_1} \ge 4$ and $|3 - 2^{k_1}| = 2^{k_1} - 3$.
In both cases $|3 - 2^{k_1}| \le 2^{k_1}$.
Since $q < 2^{k_1} x$ and, by Lemma~7, $2^{k_1 + k_2} \le 2^{k_1} q < 2^{2k_1} x$,
\[
|E| < 2 \cdot 2^{2k_1} x^2 + 2 \cdot 2^{2k_1} x + 1 = 2^{2k_1} x^2 \left(2 + \frac{2}{x} + \frac{1}{2^{2k_1} x^2}\right).
\]
Since $x \ge 3$ and $2^{k_1} \ge 2$,
\[
2 + \frac{2}{x} + \frac{1}{2^{2k_1} x^2} \le 2 + \frac{2}{3} + \frac{1}{36} = \frac{97}{36} < 3.2,
\]
so $|E| < 3.2 \cdot 2^{2k_1} x^2$.
Dividing by $2^{3k_1}$ gives $\bigl|x^3 - 2^f\bigr| < 3.2\, x^2/2^{k_1}$, and in particular
\[
\bigl|x^3 - 2^f\bigr| \le \frac{3.2\, x^2}{2^{k_1}}.
\]
Since $f = h - 3k_1$, we have $f \equiv h \pmod{3}$.
Lemma~3 says that $3$ does not divide $h$, so $3$ does not divide $f$.

\[
\bigl|x^3 - 2^f\bigr| = \bigl|x - 2^{f/3}\bigr|\bigl(x^2 + x \cdot 2^{f/3} + 2^{2f/3}\bigr).
\]
Suppose $\bigl|x - 2^{f/3}\bigr| \ge 1$.
For every real number $z$, $\bigl|x^3 - z^3\bigr|$ increases as $z$ moves away from $x$, so
\[
\bigl|x^3 - 2^f\bigr| \ge \min\bigl(x^3 - (x-1)^3,\ (x+1)^3 - x^3\bigr) = 3x^2 - 3x + 1,
\]
because $(x+1)^3 - x^3 = 3x^2 + 3x + 1 > 3x^2 - 3x + 1$.
Since $x \ge 3$, we have $3x^2 - 3x + 1 \ge 2x^2$.
Since $2^{k_1} \ge 2$, we have $3.2\, x^2/2^{k_1} \le 1.6 x^2 < 2x^2$.
So $\bigl|x^3 - 2^f\bigr| \ge 2x^2$ and $\bigl|x^3 - 2^f\bigr| < 2x^2$, which is impossible.

Hence $\bigl|x - 2^{f/3}\bigr| < 1$, so $2^{f/3} > x - 1$.
Since $x \ge 3$, we have $x - 1 \ge 2x/3$, and therefore $2^{f/3} > 2x/3$.
Then
\[
x^2 + x \cdot 2^{f/3} + 2^{2f/3} > x^2 + \frac{2x^2}{3} + \frac{4x^2}{9} = \frac{19}{9}x^2 > 2x^2,
\]
so
\[
\bigl|x - 2^{f/3}\bigr| = \frac{\bigl|x^3 - 2^f\bigr|}{x^2 + x \cdot 2^{f/3} + 2^{2f/3}} < \frac{3.2\, x^2}{2^{k_1} \cdot 2x^2} = \frac{1.6}{2^{k_1}}.
\]
In particular $\bigl|x - 2^{f/3}\bigr| \le 1.6/2^{k_1}$.
\end{proof}

\section{Proof of Theorem 1}
\label{sec:proof}

We use the following inequality, obtained in the proof of Theorem~2.1 of Bauer and Bennett (2018) from their inequality~(5.3), Lemma~4.1 and Proposition~5.1, before they round the exponent to $2^{-1.62k}$ in~(2.1).

\begin{named}{Theorem (Bauer and Bennett (2018)).}
Let $j \ge 11000$, $r \in \{1, 2\}$, $s \in \{1, 3\}$ and $n = 3j + r$.
For every integer $p$,
\[
\left| 2^{r/3} - \frac{p}{s\, 2^j} \right|
\ge
\left| 1 - \frac{p}{s\, 2^{n/3}} \right|
>
2^{-5} \cdot 16166467234^{-(n+77)/63}.
\]
\end{named}

The factor $2^{-5}$ is a weakening of the contour bound behind their inequality~(5.3). The denominator of~(5.3) contributes $1/(1.26 \cdot 2^{14/3}) = 2^{-5.00009\ldots}$, and the contour estimate for $|P_{n_1,\delta}|$ has constant $1.2528\ldots$ in place of $1.26$, so that $1.2528\ldots \cdot 2^{14/3}$ has base-$2$ logarithm $4.9918 < 5$. The formula they display omits the factor, though it is absorbed in their final statement.
Per unit of $j$ the exponent is $\mu = 3 \log_2(16166467234)/63 \approx 1.61487$.

\begin{named}{Theorem 1.}
No odd $q \ge 3$ has a cycle of $C$ of length $3$ with starting value $1$.
\end{named}

\begin{proof}[Proof of Theorem~1]
Suppose $q \ge 2^{22001}$.
Write $a = k_1$, $h = 3m + r$ and $f = h - 3a = 3j' + r'$, where $r, r' \in \{1, 2\}$ by Lemmas~\ref{lem:cube-window} and~\ref{lem:near-x}.
Lemma~5 gives $x > \sqrt{q/2} \ge 2^{11000}$.
Since $q$ is odd, $a \ge 1$, and Lemma~\ref{lem:near-x} gives $|x - 2^{f/3}| \le 0.8$.
Also $2^{f/3} \le 2^{j'+2/3}$.
If $j' \le 10999$, then $x \le 2^{10999+2/3} + 0.8$.
Now $(8/5)^3 = 512/125 > 4$, so $2^{2/3} < 8/5$, and $2^{11000} - (8/5)\,2^{10999} = (2/5)\,2^{10999} > 0.8$.
Thus $x < 2^{11000}$, a contradiction, and $j' \ge 11000$.
Likewise $(5/4)^3 = 125/64 < 2$, so $2^{1/3} > 5/4$ and $2^{j'}/4 > 0.8$, which gives $x \ge 2^{j'+1/3} - 0.8 > 2^{j'} + 2^{j'}/4 - 0.8 > 2^{j'}$.
Hence $x \ge 2^{j'}+1$ and $q = 2^a x - 1 \ge 2^{a+j'} + 2^a - 1 \ge 2^{a+j'} + 1$.
From Lemma~\ref{lem:cube-window}, $2^h > q^3$, so $2^{h/3} > q > 2^{a+j'}$ and $m + r/3 > a + j'$.
As $r \le 2$, this forces $m \ge a + j' \ge 11000$.
\begin{enumerate}
\item \textbf{The $\delta$ window.}
Apply Theorem~(Bauer and Bennett (2018)) with $s = 3$, $p = 3q + 1$, $j = m$.
Since $2^{h/3} = q + \delta$,
\[
1 - \frac{3q+1}{3 \cdot 2^{h/3}} = \frac{3(q+\delta) - (3q+1)}{3 \cdot 2^{h/3}} = \frac{\delta - 1/3}{2^{h/3}},
\]
so $\delta - 1/3 > 2^m \cdot 2^{-5} \cdot 16166467234^{-(h+77)/63}$.
With Lemma~\ref{lem:window} and $q \ge 2^m$,
\[
2^a > 1.5 \cdot 2^{2m-5} \cdot 16166467234^{-(h+77)/63}.
\]
\item \textbf{The near miss for $x$.}
Apply Theorem~(Bauer and Bennett (2018)) with $s = 1$, $p = x$, $j = j'$.
By Lemma~\ref{lem:near-x}, $|1 - x/2^{f/3}| \le 1.6/(2^a \cdot 2^{f/3})$, so
\[
2^a < 1.6 \cdot 2^{5 - f/3} \cdot 16166467234^{(f+77)/63}.
\]
\item \textbf{Combining.}
With $h = f + 3a$, $m \ge a + j'$ and $j' \ge (f-2)/3$, the $\delta$ window gives $a > 0.20878\, f - 76.8$ and the near miss for $x$ gives $a < 0.20496\, f + 47.2$.
These are incompatible once $f \ge 32{,}461$, that is $j' \ge 10{,}820$.
\end{enumerate}
So any cycle of length $3$ with starting value $1$ has $q < 2^{22001}$.
Lemma~\ref{lem:cube-window} gives $2^h < (q+1)^3$.
Since $q$ is an integer, $q + 1 \le 2^{22001}$, so $2^h < 2^{66{,}003}$ and $h < 66{,}003$.
Each such $h$ gives a single candidate $q = \lfloor 2^{h/3} \rfloor$.
Appendix~\ref{app:computations} tests every $h \le 200{,}000$ and finds no cycle, and the digit scan there rules out every such cycle with $q < 2^{1{,}071{,}999{,}953}$.
\end{proof}

\topic{Why the golden ratio appears.}
The proof squeezes $x$ between two bounds, both governed by the exponent $\mu$ in Theorem~(Bauer and Bennett (2018)).
\begin{itemize}
\item \textbf{The $\delta$ window keeps $x$ small.}
By Lemma~\ref{lem:window}, $2^{h/3}$ exceeds $q + 1/3$ by about $1/(3x)$, so $(3q+1)/(3 \cdot 2^m)$ approximates $2^{r/3}$ with error about $1/(xq)$.
Theorem~(Bauer and Bennett (2018)) forbids errors below about $q^{-\mu}$, so $x \lesssim q^{\mu-1}$.
\item \textbf{The near miss for $x$ keeps $x$ large.}
By Lemma~\ref{lem:near-x}, $x/2^{j'}$ approximates $2^{r'/3}$ with error about $1/q$, where $2^{j'} \asymp x$.
Theorem~(Bauer and Bennett (2018)) forbids errors below about $x^{-\mu}$, so $x \gtrsim q^{1/\mu}$.
\end{itemize}
Both bounds can hold only if $1/\mu \le \mu - 1$, that is $\mu^2 - \mu - 1 \ge 0$, whose positive root is $\varphi = (1 + \sqrt{5})/2$.
With $\mu = 1.61487$ the bounds read $q^{0.6192} \lesssim x \lesssim q^{0.6149}$, which are incompatible once $q$ is large; the explicit constants in the $\delta$ window, the near miss for $x$, and the comparison of those two bounds make this precise for $q \ge 2^{22001}$.
The two bounds are tied together by $q + 1 = 2^{k_1} x$. If $x$ is large, then $2^{k_1}$ is small, and the error in $(3q+1)/(3 \cdot 2^m)$ as an approximation to $2^{r/3}$ falls below the bound forbidden by Theorem~(Bauer and Bennett (2018)). If $x$ is small, then $2^{k_1}$ is large, and the error in $x/2^{j'}$ as an approximation to $2^{r'/3}$ falls below that same bound.

The threshold $\varphi$ comes from this argument; the exponent $\mu$ comes from Bauer and Bennett (2018), who needed only an exponent below $2$ for their own purposes.
Their stated value $1.62$ lies just above $\varphi$ and would leave room for $x$; the unrounded $1.61487$ lies just below it.

\section{Longer cycles}
\label{sec:longer}

\begin{named}{Lemma 10 (between consecutive $\ell$-th powers).}
\markstmt{10}{lem:power-window}
For a cycle of length $\ell$ with starting value $1$, let $\sigma = \sum_{i=1}^{\ell-1} 1/x_i$.
Then
\[
2^h = (q+1) \prod_{i=1}^{\ell-1} \left( q + \frac{1}{x_i} \right).
\]
If $\ell \ge 2$, then $q^\ell < 2^h < (q+1)^\ell$, $\ell$ does not divide $h$, and $q = \lfloor 2^{h/\ell} \rfloor$.
Also $\delta - 1/\ell = \sigma/\ell + O(1/q)$.
\end{named}

\begin{proof}
For $1 \le i \le \ell$, $2^{k_i} x_i = q x_{i-1} + 1$.
Since $h = k_1 + \cdots + k_\ell$, $x_0 = 1$ and $x_\ell = 1$,
\[
2^h = \prod_{i=1}^{\ell} \frac{q x_{i-1} + 1}{x_i} = (q+1) \prod_{i=1}^{\ell-1} \left( q + \frac{1}{x_i} \right).
\]
When $\ell = 1$ there is no index $i$ with $1 \le i \le \ell-1$, so that product has no factors and equals $1$. Thus $2^h = q + 1$.

Suppose $\ell \ge 2$.
Each $x_i$ with $1 \le i \le \ell-1$ is an odd integer at least $3$, so $q < q + 1/x_i < q + 1$.
Therefore
\[
q^\ell < (q+1)\, q^{\ell-1} < (q+1) \prod_{i=1}^{\ell-1} \left( q + \frac{1}{x_i} \right) < (q+1)^\ell.
\]
The product $(q+1) \prod_{i=1}^{\ell-1} (q + 1/x_i)$ equals $2^h$, so $q^\ell < 2^h < (q+1)^\ell$.
A power of $2$ strictly between consecutive $\ell$-th powers is not an $\ell$-th power, so $\ell$ does not divide $h$, and $q = \lfloor 2^{h/\ell} \rfloor$.

Since $2^{h/\ell} = q + \delta$, we have $(q+\delta)^\ell = 2^h$.
Dividing by $q^\ell$ gives
\[
\left(1 + \frac{\delta}{q}\right)^\ell = \left(1 + \frac{1}{q}\right) \prod_{i=1}^{\ell-1} \left(1 + \frac{1}{q x_i}\right).
\]
Expanding the product, $\prod_{i=1}^{\ell-1} (1 + 1/(q x_i)) = 1 + \sigma/q + O(1/q^2)$: the sum of the factors $1/(q x_i)$ is $\sigma/q$, and every product of two or more such factors is $O(1/q^2)$, with the constant depending only on $\ell$.
Thus
\[
\left(1 + \frac{1}{q}\right)\left(1 + \frac{\sigma}{q} + O(1/q^2)\right) = 1 + \frac{1+\sigma}{q} + O(1/q^2).
\]
Since $0 < \delta < 1$, the binomial expansion gives $(1 + \delta/q)^\ell = 1 + \ell\delta/q + O(1/q^2)$.
From $1 + \ell\delta/q + O(1/q^2) = 1 + (1+\sigma)/q + O(1/q^2)$ we have $\ell\delta/q = (1+\sigma)/q + O(1/q^2)$, so $\delta - 1/\ell = \sigma/\ell + O(1/q)$.
\end{proof}

\begin{named}{Lemma 11 (the relative error $\rho_i$).}
\markstmt{11}{lem:near-misses}
For $1 \le i \le \ell - 1$, let
\[
\varepsilon_i = \frac{1}{q} + \sum_{u=1}^{i-1} \frac{1 + \varepsilon_u}{q x_u}.
\]
Then $x_i \, 2^{S_i} = q^i (1 + \varepsilon_i)$.
Since $n_i = ih - \ell S_i$, we have $x_i = 2^{n_i/\ell}(1 + \rho_i)$.
For $q \ge \ell$,
\[
-\frac{i}{q+1} < \rho_i \le \frac{2i+1}{3q}.
\]
In particular $|\rho_i| \le (\ell - 1)/q$.
If $q \ge \ell$ and $\ell$ divides $n_i$, then $x_i \ge q/(\ell-1) - 1$.
\end{named}

\begin{proof}
For $i = 1$, we have $\varepsilon_1 = 1/q$ and $x_1 2^{k_1} = q + 1 = q(1 + 1/q)$.

Suppose $i \ge 2$ and $x_{i-1} 2^{S_{i-1}} = q^{i-1}(1 + \varepsilon_{i-1})$.
Since $2^{k_i} x_i = q x_{i-1} + 1$ and $S_i = S_{i-1} + k_i$,
\[
x_i 2^{S_i} = 2^{S_{i-1}}(q x_{i-1} + 1) = q^i(1 + \varepsilon_{i-1}) + 2^{S_{i-1}}.
\]
The inductive equation gives $2^{S_{i-1}}/q^i = (1 + \varepsilon_{i-1})/(q x_{i-1})$.
The definition of $\varepsilon_i$ gives $\varepsilon_i = \varepsilon_{i-1} + (1 + \varepsilon_{i-1})/(q x_{i-1})$.
Therefore $x_i 2^{S_i} = q^i(1 + \varepsilon_i)$.

Since $n_i = ih - \ell S_i$ and $x_i = 2^{n_i/\ell}(1 + \rho_i)$,
\[
1 + \rho_i = \frac{x_i 2^{S_i}}{2^{ih/\ell}} = \left(\frac{q}{q+\delta}\right)^i (1 + \varepsilon_i),
\]
where $2^{h/\ell} = q + \delta$.
For $\ell \ge 2$, Lemma~\ref{lem:power-window} gives $0 < \delta < 1$.

Now $\varepsilon_1 = 1/q = (2\cdot 1 + 1)/(3q)$.
If $\varepsilon_u \le 1$ for every $u < i$, then $(1 + \varepsilon_u)/(q x_u) \le 2/(3q)$, because $x_u \ge 3$, and
\[
\varepsilon_i \le \frac{1}{q} + \frac{2}{q}\cdot\frac{i-1}{3} = \frac{2i+1}{3q}.
\]
For $q \ge \ell$ and $i \le \ell - 1$, we have $(2i+1)/(3q) \le (2\ell - 1)/(3\ell) < 1$, so $\varepsilon_i \le 1$.
Thus $\varepsilon_i \le (2i+1)/(3q)$ for every $i \le \ell - 1$.
Since $0 < \delta < 1$, we have $\bigl(q/(q+1)\bigr)^i < \bigl(q/(q+\delta)\bigr)^i$, and $\bigl(q/(q+1)\bigr)^i = \bigl(1 - 1/(q+1)\bigr)^i \ge 1 - i/(q+1)$.
Also $\bigl(q/(q+\delta)\bigr)^i \le 1$ and $1 + \varepsilon_i \le 1 + (2i+1)/(3q)$, so
\[
1 - \frac{i}{q+1} \le \left(1 - \frac{1}{q+1}\right)^i < \left(\frac{q}{q+\delta}\right)^i (1 + \varepsilon_i) \le 1 + \frac{2i+1}{3q}.
\]
Here $\bigl(q/(q+\delta)\bigr)^i (1 + \varepsilon_i) = 1 + \rho_i$, so $-i/(q+1) < \rho_i \le (2i+1)/(3q)$.
Also $(2i+1)/(3q) \le (\ell-1)/q$ and $i/(q+1) < (\ell-1)/q$, so $|\rho_i| \le (\ell-1)/q$.

Suppose $q \ge \ell$ and $\ell$ divides $n_i$, so $n_i/\ell$ is an integer.
If $n_i \ge 0$, then $2^{n_i/\ell}$ is a positive integer power of $2$, hence $1$ or even, while $x_i$ is odd and at least $3$, so $|x_i - 2^{n_i/\ell}| \ge 1$.
If $n_i < 0$, then $2^{n_i/\ell} \le 1/2$ and $x_i \ge 3$, so $|x_i - 2^{n_i/\ell}| \ge 5/2$.
In either case $|\rho_i| = |x_i - 2^{n_i/\ell}| \cdot 2^{-n_i/\ell} \ge 2^{-n_i/\ell}$.
Since $2^{n_i/\ell} = x_i/(1+\rho_i)$ and $|\rho_i| \le (\ell-1)/q < 1$,
\[
2^{-n_i/\ell} = \frac{1+\rho_i}{x_i} \ge \frac{1 - |\rho_i|}{x_i}.
\]
Hence $|\rho_i| \ge (1 - |\rho_i|)/x_i$, so $x_i \ge 1/|\rho_i| - 1 \ge q/(\ell-1) - 1$.
\end{proof}

The outline in the introduction is the case $R = \{1, \ldots, \ell-1\}$.

\begin{named}{Theorem 2.}
Let $\ell \ge 2$ and let $\varphi = (1 + \sqrt{5})/2$.
Suppose that $0 < \mu < \varphi$ and that $\mu$ is a restricted irrationality exponent for length $\ell$.
Then there is a number $Q$, computed from $\ell$, $\mu$, $c(1)$ and $j_0(1)$, such that every cycle of length $\ell$ with starting value $1$ has $q < Q$.
\end{named}

\begin{proof}[Proof of Theorem~2]
Let $R = \{1, \ldots, \ell-1\}$.
Let $c$ be the positive real number $c(1)$, and let $j_0$ be the nonnegative integer $j_0(1)$.
The finitely many integers $j$ with $0 \le j < j_0$ are covered by replacing $c$ with a smaller positive real number, so the inequality holds for every integer $j \ge 0$.
Let $c'$ be the positive real number $c\, \ell^{-\mu}$.
Suppose that every multiple of $h$ has its nonzero residues modulo $\ell$ in $R$.
Write $h = \ell m + r$ with $r \in R$, and assume $q \ge 2(\ell-1)$, so Lemma~\ref{lem:near-misses} gives $|\rho_i| \le 1/2$.
The $\delta$ window uses $2^{r/\ell}$, and $x_i$ uses $2^{n_i/\ell}$ with $n_i \equiv ir \pmod{\ell}$, so the residue modulo $\ell$ of each of these exponents lies in $R$.
The definition at $x = 1$ gives, for every integer $p$ and every integer $j \ge 0$,
\[
\left| 2^{r/\ell} - \frac{p}{\ell \cdot 2^j} \right| > c \cdot 2^{-\mu j}.
\]
The same lower bound holds for the denominator $2^j$, since $p/2^j = (\ell p)/(\ell \cdot 2^j)$.
Set
\[
A = 2^{1/\ell - 1}\Bigl(\frac{c'}{2(\ell-1)}\Bigr)^{1/\mu},
\qquad
\gamma = 1 - \mu + \frac1\mu > 0,
\]
the last because $\mu^2 - \mu - 1 < 0$.

Since $2^m = (q+\delta)/2^{r/\ell} > q/2$, the definition applies to the $\delta$ window.
\[
\Bigl|\delta - \frac1\ell\Bigr|
= 2^m \Bigl| 2^{r/\ell} - \frac{\ell q + 1}{\ell \cdot 2^m} \Bigr|
> c' \cdot 2^{(1-\mu)m}.
\]
If $\mu > 1$ then $2^m < q+1$ gives $|\delta - 1/\ell| > c'(q+1)^{1-\mu} \ge c'\, 2^{1-\mu} q^{1-\mu}$.
If $\mu \le 1$ then $2^m > q/2$ gives $|\delta - 1/\ell| > c'(q/2)^{1-\mu} \ge c'/2$.

For each $x_i$ with $\ell \nmid n_i$, write $n_i = \ell j + r'$ and $x_i/2^j = 2^{r'/\ell}(1+\rho_i)$, so $|2^{r'/\ell} - x_i/2^j| = 2^{r'/\ell}|\rho_i|$.
Also $x_i/3 \le 2^j \le 2x_i$.
The inequality for the denominator $2^j$ and $|\rho_i| \le (\ell-1)/q$ give $c' \cdot 2^{-\mu j} < 2(\ell-1)/q$, so
\[
2^j > \Bigl(\frac{c'q}{2(\ell-1)}\Bigr)^{1/\mu},
\qquad
x_i \ge 2^{1/\ell-1}\, 2^j \ge A\, q^{1/\mu}.
\]
If $\ell \mid n_i$, then $x_i \ge q/(\ell-1) - 1 \ge q/(2(\ell-1))$.
Every $x_i$ is at least $\min\bigl(A q^{1/\mu},\ q/(2(\ell-1))\bigr)$, so
\[
\sigma \le (\ell-1)\, A^{-1} q^{-1/\mu} + \frac{2(\ell-1)^2}{q}.
\]

Expanding the product in Lemma~\ref{lem:power-window},
\[
\Bigl(1 + \frac{\delta}{q}\Bigr)^\ell
= \Bigl(1 + \frac1q\Bigr)\prod_{i=1}^{\ell-1}\Bigl(1 + \frac1{qx_i}\Bigr).
\]
The product is $1 + \sigma/q + \rho$ with $0 < \rho \le (\sigma/q)^2 e^{\sigma/q}$.
Comparing with the binomial expansion of the left side, and using $\sigma \le (\ell-1)/3$ and $q \ge \ell$,
\[
\Bigl|\delta - \frac1\ell - \frac{\sigma}{\ell}\Bigr|
\le \frac{\sigma}{\ell q} + \frac{\rho(q+1)}{\ell} + \frac{\ell\, e^{\ell/q}}{q}
\le \frac{4\ell}{q}.
\]
Hence
\[
\bigl|\delta - \frac1\ell\bigr|
\le A^{-1} q^{-1/\mu} + \Biggl(\frac{2(\ell-1)^2}{\ell} + 4\ell\Biggr)\frac1q.
\]

If $\mu > 1$, then $1/q \le q^{-1/\mu}$ for every integer $q \ge 1$, so
\[
\bigl|\delta - \frac1\ell\bigr|
\le \Biggl(A^{-1} + \frac{2(\ell-1)^2}{\ell} + 4\ell\Biggr) q^{-1/\mu}.
\]
Let $Q_*$ be the positive real number
\[
Q_* = \Biggl( \frac{A^{-1} + 2(\ell-1)^2/\ell + 4\ell}{c'\, 2^{1-\mu}} \Biggr)^{1/\gamma}.
\]
Take $Q = \max\bigl(2(\ell-1),\ Q_*\bigr)$.
If $q > Q$, then $\bigl|\delta - 1/\ell\bigr| > c'\, 2^{1-\mu} q^{1-\mu}$ and $\bigl|\delta - 1/\ell\bigr| \le \bigl(A^{-1} + 2(\ell-1)^2/\ell + 4\ell\bigr) q^{-1/\mu}$ cannot both hold.

If $\mu \le 1$, then $q^{-1/\mu} \le 1/q$ for every integer $q \ge 1$, so
\[
\bigl|\delta - \frac1\ell\bigr|
\le \Biggl(A^{-1} + \frac{2(\ell-1)^2}{\ell} + 4\ell\Biggr)\frac1q.
\]
Let $Q_*$ be the positive real number
\[
Q_* = \frac{2}{c'}\Biggl(A^{-1} + \frac{2(\ell-1)^2}{\ell} + 4\ell\Biggr).
\]
Take $Q = \max\bigl(2(\ell-1),\ Q_*\bigr)$.
If $q > Q$, then $\bigl|\delta - 1/\ell\bigr| > c'/2$ and $\bigl|\delta - 1/\ell\bigr| \le \bigl(A^{-1} + 2(\ell-1)^2/\ell + 4\ell\bigr)/q$ cannot both hold.
\end{proof}

For $R = \{1, \ldots, \ell-1\}$ the conclusion is that no $q > Q$ has a cycle of length $\ell$ with starting value $1$.
For $\ell = 4$ and $R = \{2\}$, the multiples of $h$ avoid the residues $1$ and $3$ exactly when $h$ is even, so the conclusion is that $h$ is odd.

For $\ell = 3$, the inequality used in the proof of Theorem~2, at $x = 1$, is the inequality supplied by Theorem~(Bauer and Bennett (2018)) for the fractions that appear. Section~\ref{sec:proof} carries out that case with explicit constants.

\begin{named}{Corollary 12.}
No cycle of length $4$ with starting value $1$ has total number of halvings $\equiv 2 \pmod{4}$.
\end{named}

\begin{proof}
Write $h = 4m + 2$ for an integer $m \ge 0$.
Then $n_1 \equiv n_3 \equiv 2 \pmod{4}$ and $n_2 \equiv 0 \pmod{4}$.
For $i \in \{1, 3\}$, let $j$ be the integer $(n_i - 2)/4$, so $x_i = \sqrt{2}\, 2^{j}(1+\rho_i)$.
Also $2^m \sqrt{2} = q + \delta$, and
\[
\Bigl|\delta - \frac14\Bigr| = 2^m \Bigl|\sqrt{2} - \frac{4q+1}{2^{m+2}}\Bigr|.
\]
Bauer and Bennett's Corollary~1.6 (2002) gives $|\sqrt{2} - p/2^j| > 2^{-1.48 j}$ for every integer $j > 2$ with $j \notin \{3, 7, 8\}$.

Let $q \ge 10^5$.
Lemma~\ref{lem:near-misses} gives $|\rho_i| \le 3/q$.
Since $\sqrt{2} < 99/70$, we have $3\sqrt{2} < 297/70 < 4.3$, so $3\sqrt{2}/q < 4.3/q$.
If $j \le 8$ and $|\rho_i| \le 3/q$, then $|x_i - 2^j \sqrt{2}| \le 3 \cdot 10^{-5} \cdot 2^8 \cdot 99/70 < 0.011$.
Hence $|x_i^2 - 2^{2j+1}| < 0.011\bigl(2 \cdot 2^8 \cdot 99/70 + 0.011\bigr) < 8$, so $|x_i^2 - 2^{2j+1}| \le 7$.
The element $x_i$ is an odd integer at least $3$.
For any two distinct odd integers $y$ and $z$ at least $3$, we have $|y^2 - z^2| \ge 16$.
Thus, for this integer $j$, at most one odd integer $x \ge 3$ satisfies $|x^2 - 2^{2j+1}| \le 7$.
Comparing $2^{2j+1}$ with those squares for $0 \le j \le 8$ leaves the pairs $(x_i, j) = (3, 0)$, $(3, 1)$, $(5, 2)$, $(11, 3)$ and $(181, 7)$.
For $181$ and $j = 7$, we have $2^{15} - 181^2 = 7$ and $2^{15/2} < 128 \cdot 99/70 < 181.03$, so
\[
\bigl|181 \cdot 2^{-15/2} - 1\bigr| = \frac{7}{2^{15/2}(181 + 2^{15/2})} > \frac{7}{70000} = 10^{-4} > \frac{3}{q}.
\]
For the other four pairs, $j \le 3$, $|x_i - 2^j \sqrt{2}| > 1/10$ and $2^j \sqrt{2} < 12$, so $|\rho_i| > 1/120 > 3/q$.
Thus $j \ge 9$.
Also $2^m = (q+\delta)/\sqrt{2} > 70q/99 > 2^{16}$, so $m \ge 17$ and $m+2 \ge 19$.
Corollary~1.6 therefore applies to $(4q+1)/2^{m+2}$ and to $x_i/2^j$.

Since $2^m < (q+1)/\sqrt{2}$,
\[
\Bigl|\delta - \frac14\Bigr| > 2^m \cdot 2^{-1.48(m+2)} > 2^{-2.72}(q+1)^{-0.48}.
\]
Also $\rho_i > -3/(q+1)$.
Since $\sqrt{2} > 41/29$ and $q \ge 10^5$, we have $\sqrt{2}\bigl(1 - 3/(q+1)\bigr) > (41/29)(1 - 3/10^5)$.
Here $41 \cdot 99997 > 141 \cdot 29 \cdot 1000$, so $(41/29)(1 - 3/10^5) > 1.41$, and therefore $2^j < x_i/1.41$.
Combined with $\sqrt{2}\,|\rho_i| \le 3\sqrt{2}/q < 4.3/q$,
\[
\frac{4.3}{q} > \Bigl(\frac{x_i}{1.41}\Bigr)^{-1.48},
\]
so $x_i > 1.41\,(q/4.3)^{1/1.48}$.
Lemma~\ref{lem:near-misses} gives $x_2 \ge q/3 - 1$.
For $q > 93$ we have $q/3 - 1 > q/3.1$, so $1/x_2 < 3.1/q$.
With $\sigma = 1/x_1 + 1/x_2 + 1/x_3$,
\[
\frac{\sigma}{4} < \frac{4.3^{1/1.48}}{2 \cdot 1.41}\, q^{-1/1.48} + \frac{0.775}{q} < 0.951\, q^{-1/1.48} + \frac{0.775}{q},
\]
where $4.3^{1/1.48}/(2 \cdot 1.41) < 0.951$.

Each $x_i \ge 3$, so $(q+\delta)^4 = (q+1)\prod_{i=1}^{3}(q + 1/x_i) \le (q+1)(q+1/3)^3$.
The difference $(q+1/2)^4 - (q+1)(q+1/3)^3 = q^2/6 + 7q/54 + 11/432$ is positive, so $\delta < 1/2$.
Write $\prod_{i=1}^{3}(1 + 1/(q x_i)) = 1 + \sigma/q + \rho$, where $\rho$ is the sum of the products of two or three factors $1/(q x_i)$.
Then $0 < \rho \le 1/(3q^2) + 1/(27q^3)$, and Lemma~\ref{lem:power-window} gives
\[
\delta - \frac14 - \frac{\sigma}{4} = \frac{q\rho}{4}\Bigl(1+\frac1q\Bigr) + \frac{\sigma}{4q} - \frac{3\delta^2}{2q} - \frac{\delta^3}{q^2} - \frac{\delta^4}{4q^3}.
\]
Since $\sigma \le 1$ and $\delta < 1/2$, for $q \ge 10^5$ the absolute value of the right-hand side is at most
\[
\frac{1+1/q}{12q} + \frac{1+1/q}{108q^2} + \frac{1}{4q} + \frac{3}{8q} + \frac{1}{8q^2} + \frac{1}{64q^3} < \frac{0.71}{q}.
\]
Thus
\[
\Bigl|\delta - \frac14\Bigr| < 0.951\, q^{-1/1.48} + \frac{1.49}{q}.
\]
At $q = 10^5$ we have $2^{-2.72}(q+1)^{-0.48} > 6.04 \times 10^{-4}$ and $0.951\, q^{-1/1.48} + 1.49/q < 4.13 \times 10^{-4}$.
Dividing the two bounds by $(q+1)^{-0.48}$ leaves the constant $2^{-2.72}$ on the left.
Each term on the right then has the form $c(q+1)^{0.48} q^{-a}$, with a positive real number $c$ and a real number $a > 0.48$, so its logarithmic derivative $0.48/(q+1) - a/q$ is negative.
The right side decreases for $q \ge 10^5$, and the inequality persists.
Appendix~\ref{app:computations} searches every $h \le 100{,}000$, and therefore every $q < 2^{25{,}000}$.
\end{proof}

\begin{named}{Corollary 13.}
Let $\ell = 2^a$ with $a \ge 2$.
Only finitely many odd $q$ have a cycle of $C$ of length $\ell$ with starting value $1$ whose total number of halvings is $\equiv \ell/2 \pmod{\ell}$, and a bound is computable from Bauer and Bennett's Corollary~1.6 (2002).
For $\ell = 4$ every such $q$ is ruled out by Corollary~12.
\end{named}

\begin{proof}
Write $h = \ell m + \ell/2$.
Every multiple of $h$ is congruent to $0$ or $\ell/2$ modulo $\ell$, so for an integer $t$ every $2^{t/\ell}$ that occurs is an integer power of $2$ or equals $2^{1/2}$.
Let $s$ be an integer in $\{1, \ell\}$, let $p$ be an integer, and let $j \ge 9$ be an integer.
Let $v$ be the integer in $\{0, a\}$ with $v = 0$ when $s = 1$ and $v = a$ when $s = \ell$.
Since $\ell = 2^a$, we have $s \cdot 2^j = 2^{j+v}$.
Then $j+v \ge 9$, so $j+v > 2$ and $j+v \notin \{3, 7, 8\}$.
Bauer and Bennett's Corollary~1.6 (2002) therefore gives
\[
\Bigl|\sqrt{2} - \frac{p}{2^{j+v}}\Bigr| > 2^{-1.48(j+v)} = 2^{-1.48 v} \cdot 2^{-1.48 j}.
\]
Since $s \cdot 2^j = 2^{j+v}$ and $v \le a$, we have $2^{-1.48 v} \ge 2^{-1.48 a}$, and
\[
\Bigl|\sqrt{2} - \frac{p}{s \cdot 2^j}\Bigr| > 2^{-1.48 a} \cdot 2^{-1.48 j}.
\]
Let $r$ be the positive integer $\ell/2$. Then $2^{r/\ell} = 2^{1/2}$.
The inequality used in the proof of Theorem~2, at $x = 1$ and at this root, holds with $\mu = 1.48$ and with $c' = 2^{-1.48 a}$.
That proof gives the bound $Q$.
The case $\ell = 4$ is Corollary~12.
\end{proof}

\begin{named}{Corollary 14.}
Let $\ell = 3 \cdot 2^a$ with $a \ge 1$.
Only finitely many odd $q$ have a cycle of $C$ of length $\ell$ with starting value $1$ whose total number of halvings is $\equiv \ell/3$ or $2\ell/3 \pmod{\ell}$, and a bound is computable from Theorem~(Bauer and Bennett (2018)).
The length $3$ is Theorem~1, which rules out every such $q$.
\end{named}

\begin{proof}
Write $h = \ell m + r_0$ with $r_0 \in \{\ell/3,\, 2\ell/3\}$.
Every multiple of $h$ is congruent to $0$, $\ell/3$, or $2\ell/3$ modulo $\ell$, so for an integer $t$ the only irrational values of $2^{t/\ell}$ that occur are $2^{1/3}$ and $2^{2/3}$.
In the proof of Theorem~2 the fraction is $(\ell q + 1)/(\ell \cdot 2^m)$.
Its numerator is coprime to $\ell$, and $\ell \cdot 2^m = 3 \cdot 2^{m+a}$, so this is a fraction of the shape in Theorem~(Bauer and Bennett (2018)) with $s = 3$ and $j = m+a$.
Each fraction $x_i/2^j$ has odd numerator $x_i$, so Theorem~(Bauer and Bennett (2018)) applies with $s = 1$.
Both give the exponent $\mu = 3\log_2(16166467234)/63 \approx 1.61487 < \varphi$ once $j$ in Theorem~(Bauer and Bennett (2018)) is at least $11000$.
The shift by $a$, and the constant factors in Theorem~(Bauer and Bennett (2018)), contribute only a constant depending on $\ell$.
The inequality used in the proof of Theorem~2, at $x = 1$ and at these roots, holds with this $\mu$ and with a positive real number $c'$ read from that theorem.
That proof gives the bound $Q$.
The length $3$ is Theorem~1.
\end{proof}

\section{Ineffective finiteness}
\label{sec:ineffective}

\begin{named}{Ridout's theorem (1958), $S$-unit form.}
Let $\alpha$ be an algebraic irrational, let $S$ be a finite set of primes, and let $\varepsilon > 0$.
Only finitely many reduced fractions $p/Q$ with $Q > 0$ have every prime factor of $Q$ in $S$ and satisfy
\[
\left| \alpha - \frac{p}{Q} \right| < Q^{-1-\varepsilon}.
\]
\end{named}

This form gives no bound on the exceptional fractions.
Since $\alpha$ is irrational, the finitely many fractions that meet it for a constant $\alpha$ lie at a positive distance from $\alpha$, and these distances have a positive minimum.
This is why Theorem~3, which uses this form in place of a restricted irrationality exponent, gives no bound on $q$.

\begin{named}{Theorem 3 (finiteness for starting value $1$).}
For each $\ell \ge 2$, only finitely many odd $q$ have a cycle of $C$ of length $\ell$ with starting value $1$.
\end{named}

The argument below uses only the $S$-unit form of Ridout's theorem and the proof of Theorem~2.
The constants supplied by Ridout's theorem are not explicit, so the resulting bound on $q$ is not computed.

\begin{proof}
Let $\varepsilon$ be a constant positive real number such that $1 + \varepsilon < \varphi$.
Let $S$ be the finite set consisting of the prime $2$ and the prime factors of $\ell$.
Let $c$ be the constant positive real number $\ell^{-1-\varepsilon}$.
For each integer $r$ with $1 \le r \le \ell - 1$, the real number $2^{r/\ell}$ is an algebraic irrational.
Ridout's theorem, applied to $2^{r/\ell}$ with this $S$ and this $\varepsilon$, allows only finitely many reduced fractions $p/Q$, with $Q > 0$ and with every prime factor of $Q$ in $S$, such that
\[
\left| 2^{r/\ell} - \frac{p}{Q} \right| < Q^{-1-\varepsilon}.
\]
The same theorem with $\varepsilon/2$ in place of $\varepsilon$ allows only finitely many reduced fractions with
\[
\left| 2^{r/\ell} - \frac{p}{Q} \right| < Q^{-1-\varepsilon/2}.
\]
Every reduced fraction satisfies $\bigl|2^{r/\ell} - p/Q\bigr| > 0$.
If $Q > 1$ and $\bigl|2^{r/\ell} - p/Q\bigr| \le Q^{-1-\varepsilon}$, then $Q^{-1-\varepsilon} < Q^{-1-\varepsilon/2}$, so the fraction is one of those given by Ridout's theorem with $\varepsilon/2$.
The reduced fractions with $Q = 1$ are integers $p$.
Since $1 < 2^{r/\ell} < 2$, the inequality $\bigl|2^{r/\ell} - p\bigr| \le 1$ holds for at most the integers $p = 1$ and $p = 2$.
Thus only finitely many reduced fractions $p_1/Q_1$ satisfy $\bigl|2^{r/\ell} - p_1/Q_1\bigr| \le Q_1^{-1-\varepsilon}$.

Let $p/(s \cdot 2^j)$ be a fraction with an integer $p$, with $s \in \{1, \ell\}$, and with an integer $j \ge 0$.
Let $p_1/Q_1$ be that fraction written in lowest terms, so $Q_1$ is a positive integer, $Q_1$ divides $s \cdot 2^j$, and every prime factor of $Q_1$ lies in $S$.
If $\bigl|2^{r/\ell} - p_1/Q_1\bigr| > Q_1^{-1-\varepsilon}$, then $Q_1 \le s \cdot 2^j$ gives
\[
\left| 2^{r/\ell} - \frac{p}{s \cdot 2^j} \right|
> Q_1^{-1-\varepsilon}
\ge (s \cdot 2^j)^{-1-\varepsilon}
= s^{-1-\varepsilon}\, 2^{-(1+\varepsilon)j}
\ge c \cdot 2^{-(1+\varepsilon)j}.
\]
The fractions $p_1/Q_1$ that remain, those with $\bigl|2^{r/\ell} - p_1/Q_1\bigr| \le Q_1^{-1-\varepsilon}$, are finite in number for each $r$, and there are finitely many integers $r$ with $1 \le r \le \ell - 1$.
For each such fraction and each $s \in \{1, \ell\}$, the inequality $(s \cdot 2^j)^{-1-\varepsilon} < \bigl|2^{r/\ell} - p_1/Q_1\bigr|$ holds for every sufficiently large integer $j$.
Hence there exists a constant integer $j_0 \ge 1$ such that, for every integer $r$ with $1 \le r \le \ell - 1$, every integer $p$, every $s \in \{1, \ell\}$, and every integer $j \ge j_0$,
\[
\left| 2^{r/\ell} - \frac{p}{s \cdot 2^j} \right| > c \cdot 2^{-(1+\varepsilon)j}.
\]
Ridout's theorem does not list the fractions $p_1/Q_1$, so this $j_0$ is not computed.

Let $\mu$ be the real number $1 + \varepsilon$.
Then $\mu < \varphi$.
The inequality displayed above is not the definition of a restricted irrationality exponent for length $\ell$.
It gives one positive real number $c$ and one nonnegative integer $j_0$ for the denominators $s \cdot 2^j$ with $s \in \{1, \ell\}$.
The definition requires a positive real number $c(x)$ and a nonnegative integer $j_0(x)$ for each odd positive integer $x$.
The comparisons in the proof of Theorem~2, at $x = 1$, use only this inequality.
That argument produces a real number $Q$.
It does not compute $Q$.
If $q > Q$, the same proof yields a contradiction.
Hence $q \le Q$.
\end{proof}

\emph{Theorem~3 fails for $\ell = 1$.}
A Mersenne number gives the cycle $1 \to 1$.
Mathematically, for $\ell \ge 2$ an odd element other than $1$ puts a factor strictly between $q$ and $q+1$, so $\ell$ does not divide $h$ and $2^{r/\ell}$ is irrational.
Ridout's theorem then applies.
For $\ell = 1$ that product has no factors and equals $1$, so the same identity becomes $2^h = q + 1$, and the three conclusions of Lemma~\ref{lem:power-window} fail.
Here $2^{h/\ell} = 2^h$ is an integer, and Ridout's theorem does not apply.

Cycles with any starting value are treated in the companion paper~\cite{hackman-ii}.
The same finiteness, for a cycle with any starting value, is Theorem~2 of that paper.

\section{Open problems}
\label{sec:open}

\begin{enumerate}
\item \textbf{Length $4$ with $h$ odd.}
The $\delta$ window and the near misses for $x_1$ and $x_3$ involve $2^{1/4}$ and $2^{3/4}$.
Write $h = 4m + r$ with an integer $m \ge 0$ and $r \in \{1, 3\}$. The denominator in the $\delta$ window is $4 \cdot 2^m = 2^{m+2}$.
Theorem~2 would finish the case given a restricted irrationality exponent below $\varphi$ for length $4$.
The construction of Bauer and Bennett (2002, 2018) needs a fourth power exceptionally close to a power of $2$; among odd $a \le 3 \times 10^6$ the best is $19^4 = 2^{17} - 751$, of quality $\log_2(2^{17})/\log_2 751 \approx 1.78$, against $4.4$ for the cube case ($5^3 + 3 = 2^7$) and $5.3$ for $\sqrt{2}$ ($181^2 + 7 = 2^{15}$).
A different construction is likely needed.
\item \textbf{Effective results for larger $\ell$.}
For a general length $\ell$, Theorem~2 needs $\mu$ to be a restricted irrationality exponent for that length.
Conjecture~1 in full would need the hypothesis of Theorem~2 for every $\ell$, or a bound uniform in $\ell$.
\end{enumerate}

\appendix

\section{Cycles of length $1$ and $2$ with starting value 1}
\label{app:one-two}

\phantomsection\label{prop:a1}
\begin{named}{Proposition A.1 (length $1$).}
The only cycle of length $1$ has starting value $1$ and $q = 2^h - 1$.
In particular the cycle with starting value $1$ is $1 \to 1$ precisely when $q$ is a Mersenne number.
\end{named}

\begin{proof}
A cycle of length $1$ is $C(x_0) = x_0$, so $x_0(2^h - q) = 1$.
Both factors are positive integers, so $x_0 = 1$ and $2^h - q = 1$.
Thus $q = 2^h - 1$.
Conversely, if $q$ is a Mersenne number, then $q = 2^e - 1$ for an integer $e \ge 2$, so $q \cdot 1 + 1 = 2^e$ and $C(1) = 1$.
\end{proof}

\phantomsection\label{prop:a2}
\begin{named}{Proposition A.2 (length $2$).}
The only odd $q$ that has a cycle of length $2$ with starting value $1$ is $q = 5$.
\end{named}

\begin{proof}
The cycle gives $q(q+1) = 2^{k_1}(2^{k_2} - 1)$.
Multiplying by $4$ and adding $1$,
\[
(2q+1)^2 + \bigl(2^{k_1+2} - 1\bigr) = 2^{k_1+k_2+2}.
\]
This is $X^2 + D = 2^n$ with $X = 2q+1$ and $D = 2^{k_1+2} - 1$.
By the theorem of Beukers (1980/81), stated as Theorem~1.8 of Bauer and Bennett (2002), for $D \ne 7$, $23$ and $D$ not of the form $2^k - 1$ there is at most one solution, while for $D = 2^k - 1$ with $k \ge 4$ there are exactly two, $X = 1$ and $X = 2^{k-1} - 1$.
These give $q = 0$ and $q = 2^{k_1} - 1$, the cycle $1 \to 1$ traversed twice, so $k_1 \ge 2$ gives nothing new.
For $k_1 = 1$, $D = 7$ is the Ramanujan--Nagell equation with $X \in \{1, 3, 5, 11, 181\}$, and only $X = 11$ gives an odd $q$ other than $1$: $q = 5$, with $1 \to 3 \to 1$.
\end{proof}

As a check, the argument of Lemma~B.2 below gives $2^{k_2} = qx + 1 < q^2 < 2^{2d}$ with $d$ the order of $2$ modulo $q$, so $k_2 = d$, and $q(q+1) = 2^{k_1}(2^d - 1)$; for $q = 5$, $d = 4$ and $5 \cdot 6 = 2 \cdot 15$.

\section{Earlier results on cycles of length $3$ with starting value 1}
\label{sec:earlier}

These results were found before Theorem~1 and are not needed for it.
Notation is as in Sections~\ref{sec:structure} to~\ref{sec:proof}.
Bauer and Bennett's headline inequality (2018), $|x^3 - 2^n| \ge x^{4/3}$ unless $x^3 = 2^n$ or $(x, n) = (5, 7)$, and their Corollary~4 (2018), $|x^3 - 27 \cdot 2^n| \ge 3^{5/3} x^{4/3}$ unless $x^3 = 27 \cdot 2^n$ or $x \in \{4, 5, 8, 15, 19, 38, 121\}$, are used throughout; both were also checked numerically for $n < 200$ and $n < 70$ respectively.

\topic{Factorization and the order of $2$.}
\begin{named}{Lemma B.1 (factorization).}
With $q = 2^{k_1} x - 1$,
\[
\bigl(2^{k_1} x - 1\bigr) \bigl(2^{k_1} x^2 - x + 1\bigr) = 2^{k_2} \bigl(2^{k_3} - 1\bigr).
\]
\end{named}

\begin{proof}
The factors are $q$ and $qx + 1$, and $q(qx + 1) = q \cdot 2^{k_2} y = 2^{k_2}(2^{k_3} - 1)$.
\end{proof}

Completing the square in $qx + 1 = 2^{k_2} y$ gives $(2q+1)^2 + (2^{k_1+2} - 1) = 2^{k_1+k_2+2} \cdot y$, so a cycle of length $3$ with starting value $1$ gives the equation of Proposition~A.2 with an odd cofactor $y$, tied to $q$ by $qy + 1 = 2^{k_3}$.

\begin{named}{Lemma B.2 (order of $2$).}
$k_3 = d$ or $k_3 = 2d$, where $d$ is the order of $2$ modulo $q$.
In particular $d < 3 \log_2 q$.
\end{named}

\begin{proof}
$2^{k_3} = qy + 1 \equiv 1 \pmod{q}$, so $d$ divides $k_3$.
Using $2^{k_1} \ge 2$ and $2^{k_2} \ge 2$, $y \le (q^2 + q + 2)/4$, so $2^{k_3} \le q(q^2 + q + 2)/4 + 1 < q^3$.
Also $2^d \ge q + 1 > q$, so $q^3 < 2^{3d}$, and hence $k_3 < 3d$.
\end{proof}

\begin{named}{Corollary B.3 (Wieferich primes).}
Let $F = (2^d - 1)/q$ and $k_3 = md$.
Then $mF \equiv 2^{-k_2} \pmod{q}$, so $\gcd(F, q) = 1$.
In particular no Wieferich prime $q$ ($q^2 \mid 2^{q-1} - 1$) has a cycle of length $3$ with starting value $1$.
\end{named}

\begin{proof}
Since $2^d \equiv 1 \pmod{q}$, $y = (2^{md} - 1)/q = F\bigl(1 + 2^d + \cdots + 2^{(m-1)d}\bigr) \equiv mF \pmod{q}$.
Reducing $2^{k_2} y = qx + 1 \bmod q$ gives $2^{k_2} \cdot mF \equiv 1$.
For a Wieferich prime, the lifting-the-exponent lemma of Lucas (1878) shows that $q^2$ divides $2^d - 1$, since $q^2$ divides $2^{q-1} - 1$ and $q$ does not divide $(q-1)/d$; so $q \mid F$.
\end{proof}

This is the mechanism Franco and Pomerance (1995) use for the question, posed by Crandall (1978), of when $1$ is not a term of the orbit.
Together with Lemma~\ref{lem:cube-window}, Lemma~B.2 determines everything from $q$: $k_1 = v_2(q+1)$, $h$ is the unique integer with $q^3 < 2^h < (q+1)^3$, $k_3 \in \{d, 2d\}$ and $k_2 = h - k_1 - k_3$.

\topic{Finiteness via Ridout's theorem.}
\begin{named}{Proposition B.4.}
Only finitely many odd $q$ have a cycle of length $3$ with starting value $1$.
\end{named}

\begin{proof}
By Lemmas~\ref{lem:window} and~5, $\delta - 1/3 < 2^{k_1+1}/(3q)$ with $2^{k_1} = (q+1)/x < (q+1)\sqrt{2/q}$, so $\delta - 1/3 < 0.96\, q^{-1/2}$ for $q \ge 71$.
With $h = 3m + r$, $\beta = 2^{r/3}$, $P = 3q + 1$ and $Q = 3 \cdot 2^m$ this gives $0 < \beta - P/Q < 2.9\, q^{-1/2}/Q$, and since $Q < 2.39(q+1)$ we get $|\beta - P/Q| < 4.6\, Q^{-3/2} < Q^{-5/4}$ for $Q > 450$.
The reduced fraction still has a denominator built from $2$ and $3$, so Ridout's theorem with $S = \{2, 3\}$ and $\varepsilon = 1/4$ allows only finitely many such fractions; infinitely many $q$ would give infinitely many distinct ones, since the errors tend to $0$ and $\beta$ is irrational.
\end{proof}

This is the case $\ell = 3$ of Theorem~3, and Theorem~1 supersedes it.
The power $3/2$ in $|\beta - P/Q| < 4.6\, Q^{-3/2}$ comes from the lower bounds $x > \sqrt{q/2}$ and $y > x$; with only $x, y \ge 3$ the error is of order $1/Q$ and Ridout's theorem gives nothing.
The result is ineffective, which is why Theorem~1 needs Theorem~(Bauer and Bennett (2018)) instead.

\topic{Elementary special cases.}
\begin{named}{Proposition B.5 ($k_1 = k_2 = 1$).}
No cycle of length $3$ with starting value $1$ has $k_1 = k_2 = 1$.
Hence $q \not\equiv 1 \pmod{8}$.
\end{named}

\begin{proof}
Lemma~2 gives $2^h = q^3 + q^2 + 2q + 4 = (q+1)(q^2 + 2) + 2$.
With $q + 1 = 2x$ this becomes
\[
4x^3 - 4x^2 + 3x + 1 = 2^n, \qquad n = h - 1,
\]
and $q^2 + 2 = 4x^2 - 4x + 3$ divides $2^n - 1$.
Since $q$ is odd, $q^2 + 2 \equiv 3 \pmod{8}$.
Every prime $p$ dividing $q^2 + 2$ has $-2$ as a square mod $p$, so $p \equiv 1$ or $3 \pmod{8}$, and at least one is $\equiv 3 \pmod{8}$.
For that $p$, $2$ is not a square mod $p$, so the order of $2$ modulo $p$ is even and divides $n$.
So $n$ is even, and $(X, Y) = (4x,\ 4 \cdot 2^{n/2})$ is an integral point on
\[
E \colon Y^2 = X^3 - 4X^2 + 12X + 16.
\]
Using Sage (passagemath) with proofs enabled, $E$ has rank $2$ and trivial torsion, generated by $(0, 4)$ and $(1, 5)$ with saturation index $1$, and its integral points are $X \in \{0, 1, 4, 9, 84\}$, computed by the elliptic-logarithm method. $X = 4$ gives $q = 1$; $X = 84$ gives $x = 21$ but $Y = 752 = 4 \cdot 188$, and $188$ is not a power of $2$; the others do not have the form $4x$ with $x$ odd and at least $3$.
For the congruence: if $q \equiv 1 \pmod{8}$, then $k_1 = 1$ and $x \equiv 1 \pmod{4}$, so $t > k_1$ in Lemma~4, and its proof forces $k_2 = k_1 = 1$.
\end{proof}

For $q = 41$, $q^3 + q^2 + 2q + 4$ is twice a square, and $qy + 1$ is not a power of $2$, since $41 \cdot 431 + 1 = 17{,}672 = 8 \cdot 47^2$.

\begin{named}{Proposition B.6 ($k_1 = 1$, $k_2 = 2$).}
No cycle of length $3$ with starting value $1$ has $k_1 = 1$ and $k_2 = 2$.
Hence $q \not\equiv 13 \pmod{16}$.
\end{named}

\begin{proof}
Suppose $k_1 = 1$ and $k_2 = 2$.
Lemma~2 gives $2^h = q^3 + q^2 + 2q + 8$.
Multiplying,
\[
(q+2)(q^2 - q + 4) = q^3 - q^2 + 4q + 2q^2 - 2q + 8 = q^3 + q^2 + 2q + 8,
\]
so $q+2$ divides $2^h$.
Lemma~1 gives $q \ge 3 \cdot 2^{k_1} - 1$.
With $k_1 = 1$ this is $q \ge 5$, and $q$ is odd, so $q+2$ is an odd integer at least $7$.
The only positive odd divisor of $2^h$ is $1$.

For the congruence: if $q \equiv 13 \pmod{16}$, then $q+1 \equiv 14 \pmod{16}$, so $v_2(q+1) = 1$.
The first step is $q+1 = 2^{k_1} x$ with $x$ odd, so $k_1 = 1$ and $x = (q+1)/2 \equiv 7 \pmod{8}$.
The second step is $qx + 1 = 2^{k_2} y$.
Substituting $q = 2x - 1$ gives $qx + 1 = 2x^2 - x + 1$.
Since $x \equiv 7 \pmod{8}$, we have $x \equiv -1 \pmod{8}$, and
\[
2x^2 - x + 1 \equiv 2(1) - (-1) + 1 = 4 \pmod{8}.
\]
Since $2x^2 - x + 1 \equiv 4 \pmod{8}$, let $m$ be the integer $(2x^2 - x + 1 - 4)/8$. Then $2x^2 - x + 1 = 4(2m+1)$ with $2m+1$ odd, so $v_2(qx + 1) = 2$.
Since $y$ is odd, $k_2 = 2$.
Lemma~2 then gives $2^h = q^3 + q^2 + 2q + 8 = (q+2)(q^2 - q + 4)$, and Lemma~1 gives $q \ge 5$, so $q+2$ is an odd integer at least $7$ dividing $2^h$.
Hence $q \not\equiv 13 \pmod{16}$.
\end{proof}

Both are superseded by Corollary~B.8, which rules out $k_1 = 1$ altogether.

\topic{Effective bounds on $k_1$.}
\begin{named}{Proposition B.7.}
In any cycle of length $3$ with starting value $1$, $2^{k_1} > q^{1/3}/2$.
\end{named}

\begin{proof}
Apply Bauer and Bennett's Corollary~4 (2018) with $x = 3q + 1$ and $n = h$.
If $(3q+1)^3 = 27 \cdot 2^h$, then $3$ divides $3q+1$.
Since $3q+1 \equiv 1 \pmod{3}$, we have $3 \nmid 3q+1$.
The only integers in $\{4, 5, 8, 15, 19, 38, 121\}$ congruent to $1$ modulo $3$ are $4$, $19$, and $121$, and these give $q = 1$, $q = 6$, and $q = 40$.
Here $q = 1$ fails $q \ge 3$, and $q = 6$ and $q = 40$ are even.
For an odd integer $q \ge 3$, neither case occurs.

Lemma~\ref{lem:window} gives $N = 27 \cdot 2^h - (3q+1)^3 > 0$, so
\[
N = \bigl|(3q+1)^3 - 27 \cdot 2^h\bigr| \ge 3^{5/3}(3q+1)^{4/3}.
\]
Since $3q+1 > 3q$,
\[
3^{5/3}(3q+1)^{4/3} > 3^{5/3}(3q)^{4/3} = 27\, q^{4/3}.
\]
By Lemma~\ref{lem:window}, $N < 54 \cdot 2^{k_1} q$.
Hence $54 \cdot 2^{k_1} q > 27\, q^{4/3}$, and therefore $2^{k_1} > q^{1/3}/2$.
\end{proof}

\begin{named}{Corollary B.8.}
(a) $k_1 = 1$ is impossible: it would force $q < 64$, and the direct search excludes every $q < 2{,}000{,}000$; so every possible $q$ is $\equiv 3 \pmod{4}$.
(b) Together with Lemma~5, $2^{2k_1-1} - 1 < q < 2^{3k_1+3}$, so each $k_1$ allows only an explicit finite range of $q$.
(c) Since the digit scan (Appendix~\ref{app:computations}) excludes $q < 2^{1{,}071{,}999{,}953}$, any remaining $q$ would need $k_1 \ge 357{,}333{,}317$, so $q \equiv -1 \pmod{2^{357{,}333{,}317}}$.
\end{named}

Theorem~1 rules out these $q$ altogether.

\begin{named}{Proposition B.9 (upper bound on $k_1$).}
$2^{k_1} \le 3.2\, x^{2/3}$, and hence $2^{k_1} < 2.01(q+1)^{2/5}$.
\end{named}

\begin{proof}
By Lemma~\ref{lem:near-x}, $\bigl|x^3 - 2^f\bigr| \le 3.2\, x^2/2^{k_1}$ and $\bigl|x - 2^{f/3}\bigr| \le 1.6/2^{k_1} < 1$, so $x^3 \ne 2^f$.
Bauer and Bennett's headline inequality (2018) gives $\bigl|x^3 - 2^f\bigr| \ge x^{4/3}$, and therefore $2^{3k_1} x^{4/3} \le 3.2 \cdot 2^{2k_1} x^2$.
The exception $x = 5$ forces $q \le 39$ by Lemma~5.
\end{proof}

Together with Proposition~B.7, $q^{1/3}/2 < 2^{k_1} < 2.01(q+1)^{2/5}$.

\topic{The three cases of Lemma 4.}
\begin{named}{Proposition B.10 ($t > k_1$ is impossible).}
No cycle of length $3$ with starting value $1$ has $t > k_1$.
\end{named}

\begin{proof}
Lemma~4 and its proof give $t = 2k_1$ and $k_2 = k_1$, so $x = 1 + 2^{2k_1} s$ with $s$ odd.
Proposition~B.7 gives $q < 8 \cdot 2^{3k_1}$, so $x < 8 \cdot 2^{2k_1}$ and $s \in \{1, 3, 5, 7\}$.
Then $q = s \cdot 2^{3k_1} + 2^{k_1} - 1$, and the cycle equation reads $2^h = q^3 + q^2 + 2^{k_1} q + 2^{2k_1}$.
Dividing by $2^{9k_1}$,
\[
s^3 < \frac{2^h}{2^{9k_1}} < \left(s + \frac{1}{2^{2k_1}}\right)^3 \left(1 + \frac{2}{q}\right),
\]
and for $2^{k_1} \ge 8$ the right-hand side is below the next power of $2$ after $s^3$ ($2$, $32$, $128$, $512$), so $2^h/2^{9k_1}$ lies strictly between consecutive powers of $2$.
For $2^{k_1} \le 4$, $q \le 451$, which the direct search covers.
\end{proof}
(Also checked in exact arithmetic for $k_1 < 400$.)

\begin{named}{Proposition B.11 ($t = k_1$ is impossible).}
No cycle of length $3$ with starting value $1$ has $t = k_1$.
\end{named}

\begin{proof}
Here $x = 1 + 2^{k_1} s$ with $s$ odd, so $q = 2^{2k_1} s + 2^{k_1} - 1$.
Let $g = h - 6k_1$.
Expanding the cycle equation,
\[
0 < 2^g - s^3 = \frac{R}{2^{6k_1}} \le \frac{6s^2}{2^{k_1}},
\qquad
R = 3 \cdot 2^{4k_1} s^2(2^{k_1}-1) + 3 \cdot 2^{2k_1} s(2^{k_1}-1)^2 + (2^{k_1}-1)^3 + q^2 + 2^{k_1} q + 2^{k_1 + k_2},
\]
using $q < 2 \cdot 2^{2k_1} s$ and $2^{k_2} < q$ ($2^{k_1} \ge 4$). Bauer and Bennett's headline inequality (2018) gives $2^g - s^3 \ge s^{4/3}$, so $2^{k_1} \le 6\, s^{2/3}$, while Proposition~B.7 gives $8 \cdot 2^{3k_1} > q > 2^{2k_1} s$, so $2^{k_1} > s/8$.
Hence $s < 48^3$ and $2^{k_1} \le 13{,}824$, so $q < 2 \cdot 2^{2k_1} s < 5 \times 10^{13}$, which the computations cover.
The exception $s = 5$ gives $q < 4 \times 10^6$ through Proposition~B.9.
\end{proof}

\begin{named}{Proposition B.12 (a constraint when $t < k_1$).}
Suppose $t < k_1$, so that $k_2 = t$, and write $x = 2^{k_1} u - (2^{k_2}-1)$ with $u$ odd.
Then $2^{k_1}/2^{k_2} \le 3u^{2/3}$.
\end{named}

\begin{proof}
Put $\delta' = 2^{k_1 + k_2} - 2^{k_1} + 1$, so $q = 2^{2k_1} u - \delta'$, and $g = h - 6k_1$.
Expanding,
\[
0 < u^3 - 2^g
= \frac{3 \cdot 2^{4k_1} u^2 \delta' - 3 \cdot 2^{2k_1} u \delta'^2 + \delta'^3 - q^2 - 2^{k_1} q - 2^{k_1 + k_2}}{2^{6k_1}}
\le \frac{3u^2 \cdot 2^{k_2}}{2^{k_1}},
\]
since the terms after the first are negative in total and $\delta' < 2^{k_1 + k_2}$.
Bauer and Bennett's headline inequality (2018) gives $u^3 - 2^g \ge u^{4/3}$.
The exceptions are $u = 1$, where $x \le 2^{k_1}$ and Proposition~B.9 forces $2^{k_1} \le 32$, and $u = 5$, where it forces $q < 4 \times 10^6$.
\end{proof}

The positivity and upper bounds in Propositions~B.11 and~B.12 were also checked in exact arithmetic on hundreds of thousands of random parameter choices.
Propositions~B.10 and~B.11 leave only the case $t < k_1$; Theorem~1 makes the case analysis unnecessary.

\topic{Small elements and an earlier bound on $y$.}
\begin{named}{Remark (small elements).}
The elements $x$ and $y$ are odd and different from $1$, so $x, y \ge 3$, and Lemma~7 gives $y \ge 5$.
If $x = 3$, Lemma~5 gives $2^{k_1} \le 2(x-1) = 4$, so $q = 3 \cdot 2^{k_1} - 1 \in \{5, 11\}$, and neither has a cycle of length $3$ with starting value $1$.
Hence $x \ge 5$ and $y \ge 7$.
None of the proofs in the main text need more than $x \ge 3$ and $y > x$.
\end{named}

Before Lemma~6 was found, the case $y < x$ could not be excluded, and the following weaker bound was used instead.
Lemma~7 shows that $y < x$ never happens, so the second alternative below is vacuous.

\begin{named}{Proposition B.13.}
In a cycle of length $3$ with starting value $1$, either $y \ge x$ or
\[
y > \left( \frac{3q^2}{q+1} \right)^{1/3} \approx (3q)^{1/3}.
\]
Moreover, in the second case $v_2(y-1) = k_1$ exactly.
\end{named}

\begin{proof}
Multiplying $2^{k_2} y = qx + 1$ by $2^{k_1} y^2$ and using $qy = 2^{k_3} - 1$ and $q + 1 = 2^{k_1} x$ gives the exact identity
\[
2^{2k_3} + (y-2)\, 2^{k_3} + C = 2^{k_1+k_2} y^3,
\qquad
C = 2^{k_1} y^2 - y + 1.
\]
If $2^{k_3}$ divides $C$, then $C \ge 2^{k_3} = qy + 1$, which simplifies to $2^{k_1} y \ge q + 1$, that is $y \ge x$.
Otherwise the first two terms are divisible by a higher power of $2$ than $C$, so $v_2(C) = k_1 + k_2$ and $2^{k_1+k_2} \le C < 2^{k_1} y^2$.
Since the left side exceeds $2^{2k_3} > q^2 y^2$, this gives $q^2 y^2 < 2^{k_1+k_2} y^3 < 2^{k_1} y^5$, so $q^2 < 2^{k_1} y^3 \le (q+1) y^3 / 3$.
For the last claim, the first step $q + 1 = 2^{k_1} x$ gives $q \equiv -1 \pmod{2^{k_1}}$.
Lemma~1 gives $k_3 \ge k_1$, so $qy + 1 = 2^{k_3} \equiv 0 \pmod{2^{k_1}}$ and $y \equiv 1 \pmod{2^{k_1}}$.
If $v_2(y-1) > k_1$, then $v_2(C) = k_1$, forcing $k_2 = 0$, which is impossible.
\end{proof}

For $q \ge 71$ the cube root is the smaller of the two bounds, so this gave $y > 1.436\, q^{1/3}$ unconditionally, and for $q > 2{,}000{,}000$ it gave $y \ge 183$.
The identity was also checked by computer algebra: substituting the cycle relations into both sides and expanding gives the same expression.

\section{Work on cycles of length $4$ with starting value 1}
\label{app:four}

\topic{The equations of the cycle.}
Write a cycle of length $4$ with starting value $1$ as $1 \to x \to y \to z \to 1$, with
\[
q + 1 = 2^{k_1} x, \qquad qx + 1 = 2^{k_2} y, \qquad qy + 1 = 2^{k_3} z, \qquad qz + 1 = 2^{k_4}.
\]
The next three lemmas are Lemma~1 and Lemma~5 for a cycle of length $4$ with starting value $1$, together with the gain from a larger $x$ in $qz + 1 = 2^{k_4}$. They constrain $k_4$ and $x$, not $y$ or $z$.

\begin{named}{Lemma C.1 (lower bound for $k_4$).}
$k_4 \ge k_1 + 3$, and $z \equiv 1 \pmod{2^{k_1}}$.
\end{named}

\begin{proof}
The odd elements other than $1$ are at least $3$, so $x \ge 3$ and $z \ge 3$.
Since $q + 1 = 2^{k_1} x$, we have $q + 1 \ge 3 \cdot 2^{k_1}$, that is $q \ge 3 \cdot 2^{k_1} - 1$.
Since $2^{k_4} = qz + 1$,
\[
2^{k_4} \ge (3 \cdot 2^{k_1} - 1) \cdot 3 + 1 = 9 \cdot 2^{k_1} - 2.
\]
Since $2^{k_1} \ge 2$, we have $9 \cdot 2^{k_1} - 2 \ge 8 \cdot 2^{k_1} = 2^{k_1 + 3}$, and hence $k_4 \ge k_1 + 3$.
From $q + 1 = 2^{k_1} x$ we also have $q \equiv -1 \pmod{2^{k_1}}$.
Thus $qz + 1 \equiv 1 - z \pmod{2^{k_1}}$.
But $qz + 1 = 2^{k_4}$ is divisible by $2^{k_1}$, because $k_4 \ge k_1$, so $1 - z \equiv 0 \pmod{2^{k_1}}$, that is $z \equiv 1 \pmod{2^{k_1}}$.
\end{proof}

\begin{named}{Lemma C.2.}
$2^{k_1} \le 2(x-1)$ and $x > \sqrt{q/2}$.
\end{named}

\begin{proof}
By Lemma~C.1, $z \equiv 1 \pmod{2^{k_1}}$, so $qy + 1 = 2^{k_3} z$ gives $y \equiv 1 - 2^{k_3} \pmod{2^{k_1}}$.
Let $t = v_2(x-1)$ and $s = (x-1)/2^t$.
If $t \ge k_1$, then $2^{k_1} \le x - 1$.
Otherwise $k_2 = t$ and $y = 2^{k_1-t}x^2 - s \equiv 2^{k_1-t} - s \pmod{2^{k_1}}$, since $x^2 \equiv 1 \pmod{2^t}$.
So $M = 2^{k_1-t} + 2^{k_3} - (s+1)$ is a multiple of $2^{k_1}$.
If $M = 0$, then $2^{k_1-t} < s + 1$, so $2^{k_1} < 2^t(s+1) \le 2(x-1)$.
If $M \le -2^{k_1}$, then $s + 1 > 2^{k_1}$, so $2^{k_1} \le x$.
If $M \ge 2^{k_1}$, then $2^{k_3} > 2^{k_1-1}$, so $k_3 \ge k_1$ and $2^{k_1-t} \equiv s + 1 \pmod{2^{k_1}}$, giving $2^{k_1-t} = s + 1$ or $s + 1 \ge 2^{k_1}$; either way $2^{k_1} \le 2(x-1)$.
Since $q + 1 = 2^{k_1} x$,
\[
q + 1 \le 2x(x-1) = 2x^2 - 2x.
\]
Hence $q < 2x^2$, that is $x > \sqrt{q/2}$.
\end{proof}

\begin{named}{Lemma C.3.}
$k_4 \ge 2k_1 + 1$ and $k_4 \ge 3k_1 - 1$.
\end{named}

\begin{proof}
By Lemma~C.2, $x \ge 2^{k_1-1} + 1$.
Since $z \ge 3$, we have $qz + 1 = 2^{k_4} \ge 3q + 1$.
Since $q = 2^{k_1} x - 1$, we have $3q + 1 = 3 \cdot 2^{k_1} x - 2$, and therefore
\[
2^{k_4} \ge 3 \cdot 2^{2k_1-1} + 3 \cdot 2^{k_1} - 2 > 2^{2k_1},
\]
so $k_4 \ge 2k_1 + 1$.

By Lemma~C.1, $z \equiv 1 \pmod{2^{k_1}}$ and $z \ne 1$, so $z \ge 1 + 2^{k_1}$.
Then $2^{k_4} = qz + 1 \ge \bigl(2^{k_1} x - 1\bigr)\bigl(2^{k_1} + 1\bigr) + 1 = 2^{k_1}\bigl(x(2^{k_1}+1) - 1\bigr)$.
Substituting $x \ge 2^{k_1-1} + 1$ yields
\[
x(2^{k_1}+1) - 1 \ge \bigl(2^{k_1-1} + 1\bigr)\bigl(2^{k_1} + 1\bigr) - 1 = 2^{2k_1-1} + 2^{k_1-1} + 2^{k_1},
\]
hence
\[
2^{k_4} \ge 2^{3k_1-1} + 2^{2k_1-1} + 2^{2k_1} = 2^{3k_1-1} + 3 \cdot 2^{2k_1-1} > 2^{3k_1-1},
\]
so $k_4 \ge 3k_1 - 1$.
\end{proof}

\topic{The middle elements.}

\begin{named}{Lemma C.4.}
If $k_2 \le k_3$, then either $y > q$, or $y > x$ and $y^2 > qx$.
In particular $y \gg q^{3/4}$.
If $k_3 \le k_4$, then either $z > q$, or $z > y$ and $z^2 > qy$.
\end{named}

\begin{proof}
Apply Lemma~6 with $i = 2$. Then $x_1 = x$, $x_2 = y$, and $x_3 = z$.
If $k_2 \le k_3$ and $y < x$, then $y > q$; if $y > x$, then $y^2 > qx$.
In the first case $y \gg q^{3/4}$.
In the second, Lemma~C.2 gives $x > \sqrt{q/2}$, so $y^2 > q^{3/2}/\sqrt{2}$ and again $y \gg q^{3/4}$.
Apply Lemma~6 with $i = 3$. Then $x_2 = y$, $x_3 = z$, and $x_4 = x_0 = 1$.
If $k_3 \le k_4$ and $z < y$, then $z > q$; if $z > y$, then $z^2 > qy$.
\end{proof}

Lemma~C.4 bounds $y$ only when $k_2 \le k_3$, and $z$ only when $k_3 \le k_4$.
It leaves open $k_3 < k_2$ with $y < x$, where $y$ could be small, and $k_4 < k_3$, where $z$ could be small.
Those two lower bounds are the missing inputs for Proposition~D.4.
Theorem~2 of the companion paper~\cite{hackman-ii} does not need them.

\section{Work on a cycle of any length with starting value 1}
\label{app:any}

\topic{The structure of a cycle.}
The next three lemmas are Lemmas~1 to~3 for an arbitrary length.

\begin{named}{Lemma D.1 (lower bound for $k_\ell$).}
Let $\ell \ge 2$.
On a cycle of length $\ell$ with starting value $1$, $k_\ell \ge k_1 + 3$ and $x_{\ell-1} \equiv 1 \pmod{2^{k_1}}$.
\end{named}

\begin{proof}
Each of $x_1$ and $x_{\ell-1}$ is an odd element other than $1$, so $x_1 \ge 3$ and $x_{\ell-1} \ge 3$.
Since $q + 1 = 2^{k_1} x_1$, we have $q + 1 \ge 3 \cdot 2^{k_1}$ and $q \ge 3 \cdot 2^{k_1} - 1$.
Also $2^{k_\ell} = q x_{\ell-1} + 1$.
Since $q \ge 3 \cdot 2^{k_1} - 1$ and $x_{\ell-1} \ge 3$,
\[
2^{k_\ell} = q x_{\ell-1} + 1 \ge \bigl(3 \cdot 2^{k_1} - 1\bigr) \cdot 3 + 1 = 9 \cdot 2^{k_1} - 2.
\]
Since $q$ is odd, $q + 1$ is even, so $k_1 \ge 1$ and $2^{k_1} \ge 2$.
Then $9 \cdot 2^{k_1} - 2 - 8 \cdot 2^{k_1} = 2^{k_1} - 2 \ge 0$, and therefore $9 \cdot 2^{k_1} - 2 \ge 8 \cdot 2^{k_1} = 2^{k_1 + 3}$.
Hence $2^{k_\ell} \ge 2^{k_1 + 3}$, so $k_\ell \ge k_1 + 3$.

From $q + 1 = 2^{k_1} x_1$, we have $q \equiv -1 \pmod{2^{k_1}}$.
Since $k_\ell \ge k_1 + 3$, we have $k_\ell \ge k_1$, so $2^{k_\ell} \equiv 0 \pmod{2^{k_1}}$.
Since $q x_{\ell-1} + 1 = 2^{k_\ell}$,
\[
1 - x_{\ell-1} \equiv q x_{\ell-1} + 1 \equiv 2^{k_\ell} \equiv 0 \pmod{2^{k_1}}.
\]
Hence $x_{\ell-1} \equiv 1 \pmod{2^{k_1}}$.
\end{proof}

The inequality $k_\ell \ge k_1 + 3$ is Lemma~1 when $\ell = 3$, and it is the first conclusion of Lemma~C.1 when $\ell = 4$.

\begin{named}{Lemma D.2 (cycle equation).}
For a cycle of length $\ell$ with starting value $1$, set $S_0 = 0$. Then
\[
2^h = q^\ell + \sum_{j=0}^{\ell-1} 2^{S_j} q^{\ell-1-j}.
\]
\end{named}

\begin{proof}
For $1 \le m \le \ell$,
\[
2^{S_m} x_m = q^m + \sum_{j=0}^{m-1} 2^{S_j} q^{m-1-j}.
\]
The case $m = 1$ is $2^{k_1} x_1 = q + 1$.
Suppose $m \ge 2$ and
\[
2^{S_{m-1}} x_{m-1} = q^{m-1} + \sum_{j=0}^{m-2} 2^{S_j} q^{m-2-j}.
\]
Since $2^{k_m} x_m = q x_{m-1} + 1$ and $S_m = S_{m-1} + k_m$,
\[
2^{S_m} x_m = 2^{S_{m-1}}(q x_{m-1} + 1) = q \cdot 2^{S_{m-1}} x_{m-1} + 2^{S_{m-1}}.
\]
Substituting the equation for $m-1$ gives
\[
2^{S_m} x_m = q\left(q^{m-1} + \sum_{j=0}^{m-2} 2^{S_j} q^{m-2-j}\right) + 2^{S_{m-1}} = q^m + \sum_{j=0}^{m-1} 2^{S_j} q^{m-1-j}.
\]
At $m = \ell$, we have $x_\ell = 1$ and $S_\ell = h$, so the left side is $2^h$ and the right side is $q^\ell + \sum_{j=0}^{\ell-1} 2^{S_j} q^{\ell-1-j}$.
\end{proof}

Lemma~2 is the case $\ell = 3$.

\begin{named}{Lemma D.3.}
For a cycle of length $\ell$ with starting value $1$ and $\ell \ge 2$,
\[
q^\ell < 2^h < (q+1)^\ell.
\]
Hence $\ell$ does not divide $h$ and $q = \lfloor 2^{h/\ell} \rfloor$.
\end{named}

\begin{proof}
Every term $2^{S_j} q^{\ell-1-j}$ in Lemma~D.2 is positive, so $2^h > q^\ell$.
Multiplying $2^{k_i} = (q x_{i-1} + 1)/x_i$ over $1 \le i \le \ell$ gives
\[
2^h = \prod_{i=1}^{\ell} \frac{q x_{i-1} + 1}{x_i} = \frac{(q+1) \prod_{i=1}^{\ell-1} (q x_i + 1)}{\bigl(\prod_{i=1}^{\ell-1} x_i\bigr) \, x_\ell} = (q+1) \prod_{i=1}^{\ell-1} \left( q + \frac{1}{x_i} \right),
\]
since $x_\ell = 1$.
Each $x_i$ with $1 \le i \le \ell-1$ is an odd element other than $1$, so $x_i \ge 3$ and $q + 1/x_i \le q + 1/3 < q+1$.
Therefore $2^h < (q+1)^\ell$.
A power of $2$ strictly between consecutive $\ell$-th powers is not an $\ell$-th power, so $\ell \nmid h$.
Thus $q < 2^{h/\ell} < q+1$, and $q = \lfloor 2^{h/\ell} \rfloor$.
\end{proof}

Lemma~3 is the case $\ell = 3$. Lemma~\ref{lem:power-window} includes these bounds.

\topic{Reduction to lower bounds.}
Before Theorem~3, finiteness for a cycle of length $\ell$ with starting value~$1$ was reduced to the hypothesis that every element other than $1$ is at least $c \cdot q^\theta$ for constant real numbers $\theta > 0$ and $c > 0$.

\begin{named}{Proposition D.4 (reduction to lower bounds).}
Let $\ell \ge 2$.
Suppose there are constant real numbers $\theta > 0$ and $c > 0$ such that, on every cycle of length $\ell$ with starting value $1$, every element other than $1$ is at least $c \cdot q^\theta$.
Then only finitely many $q$ have a cycle of length $\ell$ with starting value $1$.
\end{named}

\begin{proof}
On such a cycle, Lemma~\ref{lem:power-window} gives $\ell \nmid h$.
Write $h = \ell m + r$ with an integer $m \ge 0$ and an integer $r$ satisfying $1 \le r \le \ell-1$.
Set $\beta = 2^{r/\ell}$, $P = \ell q + 1$ and $Q = \ell \cdot 2^m$.
Then $2^m \beta = q + \delta$, so
\[
\left|\beta - \frac{P}{Q}\right| = \frac{|\delta - 1/\ell|}{2^m}.
\]
Lemma~\ref{lem:power-window} gives $\delta - 1/\ell = \sigma/\ell + O(1/q)$.
Thus there exist a constant real number $C > 0$ and a constant integer $q_0 \ge 3$ such that $|\delta - 1/\ell - \sigma/\ell| \le C/q$ for every $q \ge q_0$.
For every integer $i$ with $1 \le i \le \ell-1$, $x_i \ge c \cdot q^\theta$, so $\sigma \le (\ell-1) c^{-1} q^{-\theta}$.
Let $\eta = \min(\theta, 1)$. Then $\eta$ is a constant real number and $\eta > 0$.
For $q \ge 1$, both $q^{-\theta} \le q^{-\eta}$ and $q^{-1} \le q^{-\eta}$, so
\[
\left|\delta - \frac1\ell\right| \le \left(\frac{\ell-1}{\ell c} + C\right) q^{-\eta}
\]
whenever $q \ge q_0$.
Since $2^{1/\ell} \le \beta < 2$ and $0 < \delta < 1$, we have $q/2 < 2^m < q+1$.
Hence $Q = \ell \cdot 2^m > \ell q/2$. For $q \ge 1$ we have $q+1 \le 2q$, so $Q < 2\ell q$ and $q^{-1-\eta} < (2\ell)^{1+\eta} Q^{-1-\eta}$.
Let $K = 2\bigl((\ell-1)/(\ell c) + C\bigr) (2\ell)^{1+\eta}$. Then $K$ is a constant real number, and
\[
\left|\beta - \frac{P}{Q}\right| < K\, Q^{-1-\eta}
\]
for every $q \ge q_0$.
Let $p/Q_1$ be the fraction $P/Q$ written in lowest terms, with an integer $Q_1 > 0$.
Every prime factor of $Q_1$ lies in the finite set $S$ consisting of $2$ and the prime factors of $\ell$.
Since $Q_1 \le Q$ and $-1-\eta < 0$,
\[
\left|\beta - \frac{p}{Q_1}\right| < K\, Q_1^{-1-\eta}.
\]
Let $\varepsilon = \eta/2$. Then $\varepsilon$ is a constant real number and $\varepsilon > 0$.
For every $Q_1 > K^{1/\varepsilon}$,
\[
\left|\beta - \frac{p}{Q_1}\right| < Q_1^{-1-\varepsilon}.
\]
Ridout's theorem, with this $S$ and this $\varepsilon$, allows only finitely many fractions $p/Q_1$ satisfying $\bigl|\beta - p/Q_1\bigr| < Q_1^{-1-\varepsilon}$.
The number $\beta$ is irrational, so each of those fractions satisfies $\bigl|\beta - p/Q_1\bigr| > 0$.
As $q \to \infty$, we have $|\beta - P/Q| < K\, Q^{-1-\eta} \to 0$, so each equals $P/Q$ for only finitely many $q$.
If instead $Q_1 \le K^{1/\varepsilon}$, then $Q_1$ is a positive integer at most $K^{1/\varepsilon}$, so only finitely many such $Q_1$ occur.
For each of them, only finitely many integers $p$ satisfy $\bigl|\beta - p/Q_1\bigr| < 1$, and $\beta$ irrational gives $\bigl|\beta - p/Q_1\bigr| > 0$.
Let $\gamma$ be the minimum of those finitely many positive real numbers $\bigl|\beta - p/Q_1\bigr|$. Then $\gamma$ is a positive real number.
For every $q$ large enough that $K\, Q^{-1-\eta} < \gamma$, we have $\bigl|\beta - P/Q\bigr| < \gamma$, so $P/Q \ne p/Q_1$ for every such pair.
Every odd integer $q < q_0$ is one of finitely many integers.
Hence only finitely many $q$ have a cycle of length $\ell$ with starting value $1$.
\end{proof}

This is an earlier and weaker result than Theorem~3. It is conditional on the hypothesis that every element other than $1$ is at least $c \cdot q^\theta$ for constant real numbers $\theta > 0$ and $c > 0$, whereas Theorem~3 is not.

\section{Computations}
\label{app:computations}

For a cycle of length $\ell$ with starting value $1$ and $\ell \ge 2$, Lemma~\ref{lem:power-window} gives $q^\ell < 2^h < (q+1)^\ell$, so $\ell \nmid h$ and $q = \lfloor 2^{h/\ell}\rfloor$.
Length $1$ is Proposition~A.1, and is not searched.
A search of every $h \le H$ with $\ell \nmid h$ therefore excludes every such cycle with $q < 2^{\lfloor H/\ell \rfloor}$: it would have $2^h < (q+1)^\ell \le 2^{\ell \lfloor H/\ell \rfloor} \le 2^H$, so $h \le H$.
Searches for a cycle with any starting value, and searches at a constant $q$, are in the companion paper~\cite{hackman-ii}.

\topic{Cycles with starting value $1$.}
For each $\ell$ and each $h \le H$ with $\ell \nmid h$, the candidate $q = \lfloor 2^{h/\ell}\rfloor$ is tested by applying $C$ to $1$, $\ell$ times.
For $\ell = 3$, every $h \le 200{,}000$ was tested.
Theorem~1 needs only $h < 66{,}003$: $q < 2^{22001}$ and Lemma~\ref{lem:cube-window} give $2^h < 2^{66{,}003}$.
For each $\ell$ from $4$ to $20$, every $h \le 100{,}000$ was tested.
No cycle with starting value $1$ was found.
\begin{center}
\begin{tabular}{ccc}
\multicolumn{3}{c}{\textbf{Cycles with starting value $1$}} \\[0.4em]
length $\ell$ & $h$ at most & no cycle with starting value $1$ for $q$ below \\
$3$ & $200{,}000$ & $2^{66{,}666}$ \\
$4$ & $100{,}000$ & $2^{25{,}000}$ \\
$5$ & $100{,}000$ & $2^{20{,}000}$ \\
$6$ & $100{,}000$ & $2^{16{,}666}$ \\
$7$ & $100{,}000$ & $2^{14{,}285}$ \\
$8$ & $100{,}000$ & $2^{12{,}500}$ \\
$9$ & $100{,}000$ & $2^{11{,}111}$ \\
$10$ & $100{,}000$ & $2^{10{,}000}$ \\
$11$ & $100{,}000$ & $2^{9{,}090}$ \\
$12$ & $100{,}000$ & $2^{8{,}333}$ \\
$13$ & $100{,}000$ & $2^{7{,}692}$ \\
$14$ & $100{,}000$ & $2^{7{,}142}$ \\
$15$ & $100{,}000$ & $2^{6{,}666}$ \\
$16$ & $100{,}000$ & $2^{6{,}250}$ \\
$17$ & $100{,}000$ & $2^{5{,}882}$ \\
$18$ & $100{,}000$ & $2^{5{,}555}$ \\
$19$ & $100{,}000$ & $2^{5{,}263}$ \\
$20$ & $100{,}000$ & $2^{5{,}000}$ \\
\end{tabular}
\end{center}

\topic{Checks of Lemmas~2, 4, 6, 8 and~9.}
The identities in these lemmas were checked by computer algebra, substituting the cycle relations and expanding both sides, or on random orbits.
Lemma~4 was also checked numerically for every $k_1 \le 7$ and every odd $x < 20{,}000$.
In the proof of Theorem~1, the two bounds on $a$ were compared for every $f$ from $33{,}000$ to $400{,}000$.

\topic{Digits of $2^{1/3}$ and $2^{2/3}$.}
A cycle of length $3$ with starting value $1$ forces places in the binary expansion of $2^{1/3}$ or of $2^{2/3}$ to alternate $0, 1, 0, 1, \ldots$, beginning with $0$.
Those places do not alternate in that way in the digits that were computed, except for $m < 100$, which the search over $h$ already covers.

\begin{named}{Proposition E.1 (digit scan).}
In a cycle of length $3$ with starting value $1$ and $q \ge 17$, write $h = 3m + r$ with $r \in \{1, 2\}$ and set $\beta = 2^{r/3}$.
Then $q = \lfloor 2^m \beta \rfloor$.
Places $m+1$ through $m + \lfloor m/2 \rfloor - 2$ in the binary expansion of the fractional part of $\beta$ alternate $0, 1, 0, 1, \ldots$, beginning with $0$.
In particular, for $m \ge 100$, places $m+1$ through $m+48$ alternate in the same way.
\end{named}

\begin{proof}
Lemma~3 gives $q = \lfloor 2^{h/3} \rfloor$.
Since $h = 3m + r$, we have $2^{h/3} = 2^m \beta$, so $q = \lfloor 2^m \beta \rfloor$.
The fractional part of $2^m \beta$ is $\delta$, and its binary digits are the digits of the fractional part of $\beta$ from place $m+1$ on.
Lemma~8 gives $0 < \delta - 1/3 < 2^{k_1+1}/(3q)$; the upper bound uses $2^{k_1+k_2} \le 2^{k_1} q$ from Lemma~7.
Since $q + 1 = 2^{k_1} x$, and Lemma~5 gives $x > \sqrt{q/2}$, we have $2^{k_1} < (q+1)\sqrt{2/q}$ and
\[
0 < \delta - \frac{1}{3} < \frac{2\sqrt{2}}{3}\left(1 + \frac{1}{q}\right) q^{-1/2}.
\]
For $q \ge 17$ we have $1 + 1/q \le 18/17$, so $\frac{2\sqrt{2}}{3}\bigl(1 + \frac{1}{q}\bigr) \le 12\sqrt{2}/17 < 1$, because $\sqrt{2} < 17/12$.
Thus $\delta - 1/3 < q^{-1/2}$.
Also $\beta > 1$, so $q = \lfloor 2^m \beta \rfloor \ge 2^m$ and $q^{-1/2} \le 2^{-m/2}$.
Let $k = \lfloor m/2 \rfloor$. Then $m/2 \ge k$, so $2^{-m/2} \le 2^{-k}$, and
\[
\frac13 < \delta < \frac13 + 2^{-k}.
\]
In binary, $1/3 = 0.010101\ldots{}_2$: place $2j-1$ is $0$ and place $2j$ is $1$.
Adding a positive amount less than $2^{-k}$ can change place $k$, and the carry stops at place $k-1$, because each $1$ in this expansion is preceded by a $0$.
Places $1$ through $k-2$ of $\delta$ therefore alternate $0, 1, 0, 1, \ldots$, beginning with $0$.
These are places $m+1$ through $m + k - 2$ of the fractional part of $\beta$.
For $m \ge 100$ we have $k - 2 \ge 48$, so places $m+1$ through $m+48$ alternate in the same way.
\end{proof}

Let $N = 1{,}072{,}000{,}000$.
The first $N$ binary places of $\beta = 2^{1/3}$ and of $\beta = 2^{2/3}$ were computed and checked by $I^3 < 2^{3N+r} < (I+1)^3$, where $I = \lfloor 2^N \beta \rfloor$.
A multiprecision value was kept only when its $192$ guard digits were neither all $0$ nor all $1$ and its leading $61$ digits matched an integer cube root.
Above about $2^{30}$ digits that value disagreed with the integer test, so $N$ was kept below $2^{30}$ and the integer test is the one used here.
For $m \ge 100$, Proposition~E.1 requires places $m+1$ through $m+48$ to alternate.
Places $m+1$ through $m+48$ do not alternate in that way for any $m$ with $m + 48 \le N$, that is $m \le N - 48$.
Such an $m$ would give $q = \lfloor 2^m \beta \rfloor < 2^{m+1} \le 2^{N-47} = 2^{1{,}071{,}999{,}953}$, since $\beta < 2$.
Every $m < 100$ has $h = 3m + r < 300$, which the search $h \le 200{,}000$ includes.
Together these rule out every cycle of length $3$ with starting value $1$ and $q < 2^{1{,}071{,}999{,}953}$.

\topic{Checking Bauer and Bennett (2018).}
The results of this paper rely heavily on Bauer and Bennett (2018), so the argument that produces their inequality was checked.
The labels below are those of the published paper.
In the preprint of Bauer and Bennett (2018) the same results are Theorem~3, inequality~(21), Lemma~7, Proposition~8 and Lemma~9.

Theorem~(Bauer and Bennett (2018)) is the inequality in the proof of their Theorem~2.1, before the exponent is rounded.
The constant $16166467234$ in that inequality is the integer just above $9 \cdot 2^{28} \cdot F(\tau)/1.79$, where $F$ is their auxiliary function: $F(\tau) = 11.97804\ldots$ at $\tau = 0.20305\ldots$, and the product equals $16166467233.3\ldots$.
The denominator of their inequality~(5.3) contributes $1/(1.26 \cdot 2^{14/3}) = 2^{-5.00009\ldots}$, and the formula they display omits this factor.
The factor $2^{-5}$ in Theorem~(Bauer and Bennett (2018)) is a weakening of the contour bound for $|P_{n_1,\delta}|$, which has constant $1.2528\ldots$ in place of $1.26$: then $1.2528\ldots \cdot 2^{14/3}$ has base-$2$ logarithm $4.9918 < 5$.
The sentence before the proof of Theorem~1 records the same comparison, and the factor is absorbed in their final bound $2^{-1.62k}$.

Their Lemma~4.1 was checked at $3000$-digit precision for every parameter $n_1$ from $2$ to $60$ and both values of each of $\delta$ and $r$.
The inequality uses $n_1 \ge 497$, and Lemma~4.1 is uniform in $n_1$.
With
\[
f_{\delta,r}(v) = \frac{v^{4+r/3-\delta}(1-v)^{1-r/3}}{\bigl(1 - 3(1-v)/128\bigr)^{5-\delta}},
\]
we have $\max f_{\delta,r} = 0.31999\ldots < 1/3$, and $F(\tau) < 12$.

Their Proposition~5.1 states that a product $\Pi$, built from their Pad\'e approximants, satisfies $\Pi > 1.8^{n_1}$ for every integer $n_1 \ge 497$.
It was recomputed as follows.
\begin{itemize}
\item Their Lemma~6.1 lists the primes that divide every Pad\'e coefficient.
For every $n_1 \le 160$ and both values of each of $\delta$ and $r$, the $5{,}978$ primes that meet its hypotheses were checked, with no failure.
\item A lower bound $L(r, n_1)$ for $\log \Pi$ was computed exactly for every $497 \le n_1 < 10^6$ and both values of $r$.
It lower-bounds $\log \Pi$ through Lemma~6.1.
The cutoff $\lfloor \sqrt{n_1}/12 \rfloor - 2$ in that application is negative for $n_1 < 576$, so on $n_1 \le 830$ the sum includes primes beyond the cutoff, each still required to satisfy $p > \max\bigl\{\sqrt{3n_2+2},\, 5\bigr\}$.
It reproduces their values $\exp(L(2, 3319)/3319) = 1.897735\ldots$ and $\exp(L(1, 10000)/10000) = 2.006523\ldots$.
The bound exceeds $\log 1.8$ except for $43$ pairs $(n_1, r)$, all with $n_1 \le 830$.
For those pairs $\Pi$ was computed exactly for both values of $\delta$, and the minimum of $\Pi^{1/n_1}$ is $1.80767$, at $n_1 = 497$.
For $1000 < n_1 < 10^6$ the minimum of $\exp(L(r, n_1)/n_1)$ is $1.8422$.
\item For $n_1 \ge 10^6$ the analytic lower estimates of $\log \Pi$ in the preprint of Bauer and Bennett (2018) were recomputed: $0.59197\, n_1$ when $10^6 \le n_1 \le 2 \times 10^9$, and $0.59035\, n_1$ when $n_1 > 2 \times 10^9$.
Both exceed $\log 1.8 \cdot n_1 = 0.58779\, n_1$.
Those estimates use the preprint's prime-error bound, in which $1.798158\sqrt{x}$ and $0.00144\, x/\log x$ meet at $x = 641{,}239{,}200.2$.
The journal states the coefficient $0.001$ and splits this argument at $3 \times 10^9$.
The constant $1.798158$, from Bennett, Martin, O'Bryant and Rechnitzer (2018), was spot-checked for $x \le 2 \times 10^8$, and the largest ratio found is $1.79628$.
\end{itemize}
What was not recomputed is that prime bound for $x > 2 \times 10^8$, Lemma~4 of Beukers (1980/81), where one choice of $\delta$ gives $\Lambda \ne 0$, and Lemma~4.5 of Chudnovsky (1983), which Lemma~6.1 uses.

\section{Notes on the literature and history}
\label{app:history}

\topic{Literature.}
\begin{itemize}
\item Santos (2020) studies cycles of $C$ with starting value $1$, the same question as this paper.
Propositions~15 and~16 there are the cases of length $1$ and length $2$, and Conjecture~3 is that a cycle of $C$ has starting value $1$ only when $q = 5$ or when $q$ is a Mersenne number.
His Remark~5 describes a cycle of length $3$ or more as ``\ldots{} much more difficult to analyze and seems intractable''; Theorem~1 is the case of length $3$.
The same preprint claims a proof of the Collatz conjecture, which is not accepted.
\item Crandall (1978) and Franco and Pomerance (1995) study when $1$ is not a term of an orbit of a $qx+1$ map; Corollary~B.3 matches their mechanism.
\item Steiner (1977, 1981) rules out one-block ``circuits'' for $q = 3$, $5$, $7$ with Baker's method, a related use of Diophantine approximation for cycles.
\item The deep inputs are Bauer and Bennett (2002, 2018) on restricted irrationality exponents, Beukers (1980/81) on $X^2 + D = 2^n$, and Ridout (1958) on approximation with $S$-unit denominators.
\end{itemize}

\topic{History of the results.}
The work began with the $\delta$ window (Lemma~\ref{lem:window}), which pins down $q$ from $h$ and makes a search over $h$ far faster than a search over $q$.
Lower bounds on $x$ and $y$ (Lemmas~4 to~7, preceded by the weaker Proposition~B.13, which allowed $y < x$ with $y > (3q)^{1/3}$) then gave finiteness through Ridout's theorem (Proposition~B.4), but no bound.
Residue-class arguments settled special cases (Propositions~B.5 and~B.6), and the effective inequality of Bauer and Bennett (2018) gave the lower bound on $k_1$ (Proposition~B.7) and most of the case analysis of Lemma~4 (Appendix~\ref{sec:earlier}). Playing the $\delta$ window against the near miss for $x$ then gave Theorem~1, once it was noticed that the unrounded exponent of Bauer and Bennett (2018) lies below the golden ratio.
Lemma~\ref{lem:near-misses} extended this to Theorem~2, Corollaries~12 to~14, and Theorem~3.
The digit scan (Appendix~\ref{app:computations}) was computational evidence found along the way.
Cycles with any starting value are treated in the companion paper~\cite{hackman-ii}.
The final verification found and corrected a dropped constant in a displayed formula of Bauer and Bennett (2018) and confirmed their Proposition~5.1 for every $n_1 \ge 497$.

\begin{thebibliography}{99}
\bibitem{bauer-bennett-2002}
M.~Bauer and M.~A.~Bennett, Applications of the hypergeometric method to the generalized Ramanujan--Nagell equation, Ramanujan J.\ 6 (2002), 209--270, \doilink{10.1023/A:1015779301077}.

\bibitem{bauer-bennett-2018}
M.~Bauer and M.~A.~Bennett, Ramanujan--Nagell cubics, Rocky Mountain J.\ Math.\ 48 (2018), 385--412, \doilink{10.1216/rmj-2018-48-2-385}.

\bibitem{bennett-martin-obryant-rechnitzer}
M.~A.~Bennett, G.~Martin, K.~O'Bryant and A.~Rechnitzer, Explicit bounds for primes in arithmetic progressions, Illinois J.\ Math.\ 62 (2018), 427--532, \doilink{10.1215/ijm/1552442669}.

\bibitem{beukers}
F.~Beukers, On the generalized Ramanujan--Nagell equation~I, Acta Arith.\ 38 (1980/81), 389--410, \doilink{10.4064/aa-38-4-389-410}.

\bibitem{chudnovsky}
G.~V.~Chudnovsky, On the method of Thue--Siegel, Ann.\ of Math.\ 117 (1983), 325--382, \doilink{10.2307/2007080}.

\bibitem{crandall}
R.~E.~Crandall, On the ``$3x+1$'' problem, Math.\ Comp.\ 32 (1978), 1281--1292, \doilink{10.1090/S0025-5718-1978-0480321-3}.

\bibitem{franco-pomerance}
Z.~Franco and C.~Pomerance, On a conjecture of Crandall concerning the $qx+1$ problem, Math.\ Comp.\ 64 (1995), 1333--1336, \doilink{10.1090/S0025-5718-1995-1297468-4}.

\bibitem{hackman-ii}
T.~Hackman, II. Cycles beyond $1$ for the $qn+1$ map: finiteness for each length, and an explicit bound on $q$ for a constant odd element, companion paper.

\bibitem{lucas}
\'E.~Lucas, Th\'eorie des fonctions num\'eriques simplement p\'eriodiques, Amer.\ J.\ Math.\ 1 (1878), 184--240, 289--321.

\bibitem{ridout}
D.~Ridout, The $p$-adic generalization of the Thue--Siegel--Roth theorem, Mathematika 5 (1958), 40--48, \doilink{10.1112/S0025579300001339}.

\bibitem{santos}
R.~Santos, On the Collatz general problem $qn+1$, preprint (2020), \href{https://arxiv.org/abs/2005.00346}{\nolinkurl{arXiv:2005.00346}}.

\bibitem{steiner-1977}
R.~P.~Steiner, A theorem on the Syracuse problem, Proc.\ 7th Manitoba Conference on Numerical Mathematics (1977), 553--559.

\bibitem{steiner-1981}
R.~P.~Steiner, On the ``$qx + 1$ problem'', $q$ odd, Fibonacci Quart.\ 19 (1981), 285--288, \doilink{10.1080/00150517.1981.12430094}.
\end{thebibliography}

\end{document}
